\documentclass[b5paper,twoside]{article}
\usepackage[fontsize=13pt]{scrextend}
\usepackage[
  paperwidth=176mm,
  paperheight=250mm,
  left=7mm,
  right=7mm,
  top=9.1mm,
  bottom=9.1mm,
  includehead,
  ignorefoot,
  headheight=18pt,
  headsep=4pt
]{geometry}

\usepackage{amsmath}
\usepackage{newpxtext}
\usepackage[amsthm]{newpxmath}
\usepackage{microtype}
\usepackage{fancyhdr}
\usepackage{xcolor}
\usepackage{hyperref}

\allowdisplaybreaks
\linespread{1.1}
\setlength{\emergencystretch}{2em}
\setlength{\parindent}{1.5em}
\setlength{\parskip}{0pt}

\fancyhf{}
\fancyhead[LE]{\thepage}
\fancyhead[RO]{\thepage}
\renewcommand{\headrulewidth}{0pt}
\renewcommand{\footrulewidth}{0pt}
\pagestyle{fancy}

\definecolor{InternalLink}{HTML}{1A4D8F}
\definecolor{CitationLink}{HTML}{2C6E49}
\definecolor{ExternalLink}{HTML}{8A2C5A}
\hypersetup{
  colorlinks=true,
  linkcolor=InternalLink,
  citecolor=CitationLink,
  urlcolor=ExternalLink,
  filecolor=ExternalLink,
  pdfborder={0 0 0},
  unicode=true,
  bookmarksopen=true,
  bookmarksnumbered=false,
  pdftitle={Grundzuege einer allgemeinen Theorie der linearen Integralgleichungen. Vierte Mitteilung / Fundamentals of a General Theory of Linear Integral Equations. Fourth Communication},
  pdfauthor={David Hilbert},
  pdfsubject={English translation of the 1906 Fourth Communication}
}

\makeatletter
\newcommand{\hilbertstatementname}[2]{%
  \textbf{#1\if\relax\detokenize{#2}\relax\else\ #2\fi.}%
}
\newenvironment{theorem}[1][]{%
  \par\addvspace{.7\baselineskip}%
  \phantomsection
  \protected@edef\@currentlabel{#1}%
  \noindent\hilbertstatementname{Theorem}{#1}\ \itshape\ignorespaces
}{%
  \par\addvspace{.7\baselineskip}%
}
\newenvironment{lemma}[1][]{%
  \par\addvspace{.7\baselineskip}%
  \phantomsection
  \protected@edef\@currentlabel{#1}%
  \noindent\hilbertstatementname{Lemma}{#1}\ \itshape\ignorespaces
}{%
  \par\addvspace{.7\baselineskip}%
}
\newenvironment{definition}[1][]{%
  \par\addvspace{.5\baselineskip}%
  \phantomsection
  \protected@edef\@currentlabel{#1}%
  \noindent\hilbertstatementname{Definition}{#1}\ \normalfont\ignorespaces
}{%
  \par\addvspace{.5\baselineskip}%
}
\renewcommand{\proofname}{\textit{Proof}}
\makeatother

\newcommand{\originalpage}[1]{\phantomsection\label{origpage:#1}}

\begin{document}
\pagenumbering{arabic}
\thispagestyle{empty}
\begin{center}
  {\LARGE\bfseries Grundz\"uge einer allgemeinen Theorie der linearen
  Integralgleichungen.\par}
  \vspace{0.35\baselineskip}
  {\Large\bfseries (Vierte Mitteilung).\par}
  \vspace{0.8\baselineskip}
  {\LARGE\bfseries Fundamentals of a General Theory of Linear Integral
  Equations\par}
  \vspace{0.35\baselineskip}
  {\Large\bfseries (Fourth Communication).\par}
  \vspace{0.8\baselineskip}
  {\large By David Hilbert.\par}
  \vspace{0.35\baselineskip}
  {Presented at the session of 3 March 1906.\par}
\end{center}
\vspace{0.6\baselineskip}

\originalpage{157}
In the introduction to my first communication I announced a new method for
the treatment of integral equations; it rests upon a theory of quadratic
forms in infinitely many variables, which I shall develop in the present
communication.

The systematic treatment of quadratic forms in infinitely many variables is
also of great importance in itself and appears to me to be an essential
completion of the familiar theory of quadratic forms in a finite number of
variables.  The applications of the theory of quadratic forms in infinitely
many variables are not confined to integral equations: this theory will come
into contact no less with Stieltjes's beautiful theory of continued
fractions\footnote{\hyperref[bib:Stieltjes1894]{\emph{Recherches sur les
fractions continues}, Ann. de Toulouse VIII (1894).}} than, on the other hand, with the
question of solving systems of infinitely many linear equations, whose study
Hill, Poincar\'e, H.~von Koch, and others have successfully undertaken.  There
is also, from the methodological point of view, a close relation to
Minkowski's investigations of volume and surface area.  Above all, however,
the theory of quadratic forms in infinitely many variables
opens\originalpage{158} a new path
to the most general expansions of arbitrary functions in infinite series of
Fourier type, as will be indicated at the end of
\hyperref[sec:XI]{Section XI} of this communication.

In a fifth communication I shall show not only how extraordinarily simply all
the results on integral equations of the second kind obtained in my first
communication can now be proved and substantially extended, but also how the
method of quadratic forms in infinitely many variables permits these results
to be \mbox{transferred} to integral equations of the first kind and to certain
integral equations of a new ``third'' kind, and how, by the new method,
integral equations whose kernel exhibits higher singularities likewise become
amenable to successful treatment.

\section*{XI. Theory of the Orthogonal Transformation of a\\
Quadratic Form in Infinitely Many Variables}
\addcontentsline{toc}{section}{XI. Theory of the Orthogonal Transformation of a Quadratic Form in Infinitely Many Variables}
\label{sec:XI}

Let
\[
  k_{pq}=k_{qp},
\]
where each of the two indices $p,q$ runs through the sequence of all positive
integers $1,2,\ldots$, be arbitrarily given real constants.  Then
\[
  K(x)=\sum_{p,q=1,2,\ldots} k_{pq}x_p x_q
\]
represents a \emph{quadratic form} in the infinitely many variables
$x_1,x_2,\ldots$; the constants $k_{pq}$ are called the coefficients of this
quadratic form.  An expression
\[
  A(x,y)=\sum_{p,q=1,2,\ldots} a_{pq}x_p y_q
\]
with arbitrary coefficients $a_{pq}$ is called a \emph{bilinear expression}
or a \emph{bilinear form} in the infinitely many variables
$x_1,x_2,\ldots,y_1,y_2,\ldots$.  The expression
\[
  K(x,y)=\sum_{p,q=1,2,\ldots} k_{pq}x_p y_q
\]
is called the bilinear form belonging to $K(x)$.  It is a form in the
infinitely many variables $x_1,x_2,\ldots$ and $y_1,y_2,\ldots$.

If we allow the indices $p,q$ to run only through the values $1,\ldots,n$,
where $n$ denotes a finite number, there arise\originalpage{159} the forms
with the finite number $n$ of variables as follows:
\[
  K_n(x)=\sum_{p,q=1,\ldots,n}k_{pq}x_p x_q,
  \qquad
  K_n(x,y)=\sum_{p,q=1,\ldots,n}k_{pq}x_p y_q.
\]
These forms $K_n(x)$ and $K_n(x,y)$ shall be called the \emph{truncations} of
the forms $K(x)$ and $K(x,y)$, respectively, and shall also be denoted by
$[K]$, sometimes without explicitly appending the index $n$.

If any bilinear expression with a finite number of variables,
\[
  A_n(x,y)=\sum_{p,q=1,\ldots,n}a_{pq}x_p y_q,
\]
is given, the sum of all the coefficients of the terms $x_p y_p$ shall be
called the \emph{convolution} of the expression $A_n(x,y)$ and shall be
denoted by $A_n(\mathord{\cdot},\mathord{\cdot})$; thus\footnote{Translator's
note: In this displayed formula the original edition omits the subscript $n$
from $A_n(\mathord{\cdot},\mathord{\cdot})$; it has been supplied here in
accordance with the preceding definition.}
\[
  A_n(\mathord{\cdot},\mathord{\cdot})
  =\sum_{p=1,\ldots,n}a_{pp},
\]
and in particular
\[
  K_n(\mathord{\cdot},\mathord{\cdot})
  =k_{11}+\cdots+k_{nn}.
\]

Besides the general quadratic forms $K,K_n$, we shall also consider the
special quadratic forms
\[
  (x,x)=x_1^2+x_2^2+\cdots,
  \qquad
  (x,x)_n=x_1^2+\cdots+x_n^2,
\]
and the special bilinear forms
\[
  (x,y)=x_1y_1+x_2y_2+\cdots,
  \qquad
  (x,y)_n=x_1y_1+\cdots+x_ny_n.
\]

The discriminant of the form
\[
  (x,x)_n-\lambda K_n(x),
\]
that is, the determinant
\begin{equation}
\tag{1}\label{eq:1}
  D_n(\lambda)=
  \begin{vmatrix}
    1-\lambda k_{11} & -\lambda k_{12} & \cdots & -\lambda k_{1n}\\
    -\lambda k_{21} & 1-\lambda k_{22} & \cdots & -\lambda k_{2n}\\
    \vdots & \vdots & \ddots & \vdots\\
    -\lambda k_{n1} & -\lambda k_{n2} & \cdots & 1-\lambda k_{nn}
  \end{vmatrix},
\end{equation}
is, as is well known, a polynomial of degree $n$ in $\lambda$, all of whose
zeros are real.  Let these be denoted by
\[
  \lambda_1^{(n)},\lambda_2^{(n)},\ldots,\lambda_n^{(n)}
\]
\originalpage{160}and called the \emph{eigenvalues} of the form $K_n$; the totality of these
$n$ eigenvalues shall be called the \emph{spectrum} of the form $K_n$.  If a
real value $\lambda$ is such that, at that value or in every neighborhood of
it, eigenvalues of $K_n$ occur for infinitely many $n$, this value $\lambda$
shall be called a \emph{condensation value} of $K$. If the greatest absolute
values among the eigenvalues of $K_n$ increase beyond every bound as $n$
increases, the value $\lambda=\infty$ shall likewise
be counted as a condensation value of the form $K$.

We now border the determinant in \eqref{eq:1}, obtaining
\[
  D_n(\lambda;x,y)=
  \begin{vmatrix}
    1-\lambda k_{11} & \cdots & -\lambda k_{1n} & x_1\\
    \vdots & & \vdots & \vdots\\
    -\lambda k_{n1} & \cdots & 1-\lambda k_{nn} & x_n\\
    y_1 & \cdots & y_n & 0
  \end{vmatrix},
\]
and call the quotient
\[
  K_n(\lambda;x,y)=-\frac{D_n(\lambda;x,y)}{D_n(\lambda)},
  \qquad
  K_n(\lambda;x)=K_n(\lambda;x,x),
\]
the \emph{resolvent of the quadratic form} $K_n$, with the second notation
being the special case $y=x$. If
$\lambda$ is not an eigenvalue, the coefficients of $x_1,\ldots,x_n$ in
$K_n(\lambda;x,y)$ give the solution of the linear equations
\[
  x_p-\lambda(k_{p1}x_1+\cdots+k_{pn}x_n)=y_p,
  \qquad (p=1,\ldots,n);
\]
we express this by the identity
\[
  K_n(\lambda;x,y)
  -\lambda K_n(x,\mathord{\cdot})K_n(\lambda;\mathord{\cdot},y)
  =(x,y)_n.
\]
It is evident that
\begin{equation}
\tag{2}\label{eq:2}
\begin{aligned}
  K_n(\lambda;x,y)
  &=(x,y)_n+\lambda K_n(x,y)+\lambda^2K_nK_n(x,y)\\
  &\quad+\lambda^3K_nK_nK_n(x,y)+\cdots,
\end{aligned}
\end{equation}
where $K_nK_n,K_nK_nK_n,\ldots$ denote the forms arising from $K_n$
by repeated convolutions,
\begin{align*}
  K_nK_n(x,y)
    &=K_n(x,\mathord{\cdot})K_n(\mathord{\cdot},y)
      =\sum_{p,q,r=1,\ldots,n}k_{pr}k_{rq}x_p y_q,\\
  K_nK_nK_n(x,y)
    &=K_n(x,\mathord{\cdot})K_nK_n(\mathord{\cdot},y)
      =\sum_{p,q,r,s=1,\ldots,n}k_{pr}k_{rs}k_{sq}x_p y_q,
\end{align*}
and so on.  For the resolvent $K_n$ the partial-fraction representation
\originalpage{161}
\begin{equation}
\tag{3}\label{eq:3}
  K_n(\lambda;x,y)
  =\frac{L_1^{(n)}(x)L_1^{(n)}(y)}{1-\lambda/\lambda_1^{(n)}}
   +\cdots+
   \frac{L_n^{(n)}(x)L_n^{(n)}(y)}{1-\lambda/\lambda_n^{(n)}}
\end{equation}
holds, where $L_1^{(n)},\ldots,L_n^{(n)}$ denote certain linear forms in
the variables with real coefficients.  In particular, therefore, the
following formulas hold:
\begin{equation}
\tag{4}\label{eq:4}
\begin{aligned}
  (x,y)_n
    &=x_1y_1+\cdots+x_ny_n
      =L_1^{(n)}(x)L_1^{(n)}(y)+\cdots+L_n^{(n)}(x)L_n^{(n)}(y),\\
  K_n(x,y)
    &=\frac{L_1^{(n)}(x)L_1^{(n)}(y)}{\lambda_1^{(n)}}
      +\cdots+
      \frac{L_n^{(n)}(x)L_n^{(n)}(y)}{\lambda_n^{(n)}};
\end{aligned}
\end{equation}
\begin{equation}
\tag{5}\label{eq:5}
\begin{aligned}
  (x,x)_n
    &=x_1^2+\cdots+x_n^2
      =(L_1^{(n)}(x))^2+\cdots+(L_n^{(n)}(x))^2,\\
  K_n(x)
    &=\frac{(L_1^{(n)}(x))^2}{\lambda_1^{(n)}}
      +\cdots+
      \frac{(L_n^{(n)}(x))^2}{\lambda_n^{(n)}},\\
  K_nK_n(x)
    &=\frac{(L_1^{(n)}(x))^2}{(\lambda_1^{(n)})^2}
      +\cdots+
      \frac{(L_n^{(n)}(x))^2}{(\lambda_n^{(n)})^2},\\
  K_nK_nK_n(x)
    &=\frac{(L_1^{(n)}(x))^2}{(\lambda_1^{(n)})^3}
      +\cdots+
      \frac{(L_n^{(n)}(x))^2}{(\lambda_n^{(n)})^3},
      \qquad\ldots
\end{aligned}
\end{equation}
From \eqref{eq:4} follow the orthogonality properties of the linear forms
$L_1^{(n)},\ldots,L_n^{(n)}$,
\[
  L_p^{(n)}(\mathord{\cdot})L_q^{(n)}(\mathord{\cdot})=0
  \quad(p\ne q),
  \qquad
  L_p^{(n)}(\mathord{\cdot})L_p^{(n)}(\mathord{\cdot})=1,
\]
where $L_p^{(n)}(\mathord{\cdot})L_q^{(n)}(\mathord{\cdot})$ denotes the
convolution of the bilinear form $L_p^{(n)}(x)L_q^{(n)}(y)$; and then,
again by convolution, one obtains
\begin{equation}
\tag{6}\label{eq:6}
\begin{aligned}
  K_n(\mathord{\cdot},\mathord{\cdot})
    &=\sum_{p=1,\ldots,n}k_{pp}
      =\frac1{\lambda_1^{(n)}}+\cdots+\frac1{\lambda_n^{(n)}},\\
  K_nK_n(\mathord{\cdot},\mathord{\cdot})
    &=\sum_{p,r=1,\ldots,n}k_{pr}k_{rp}
      =\frac1{(\lambda_1^{(n)})^2}+\cdots
       +\frac1{(\lambda_n^{(n)})^2},\\
  K_nK_nK_n(\mathord{\cdot},\mathord{\cdot})
    &=\sum_{p,r,s=1,\ldots,n}k_{pr}k_{rs}k_{sp}
      =\frac1{(\lambda_1^{(n)})^3}+\cdots
       +\frac1{(\lambda_n^{(n)})^3},
       \qquad\ldots
\end{aligned}
\end{equation}

Our first important task is to transfer the concept of the resolvent of the
form $K_n$ to the form $K$, and to seek, for the form $K$ in infinitely many
variables, the analogue of the partial-fraction representation
\eqref{eq:3}.

For this purpose we use a general lemma from the domain of continuous
(uniform) convergence,\originalpage{162} which I have already applied in a more special form in
my investigations on Dirichlet's principle.\footnote{\hyperref[bib:Hilbert1904Dirichlet]{
\emph{Math. Ann.}, vol.~59 (1904).}}

\begin{lemma}[1]\label{lem:1}
Let $J$ be an interval lying wholly in the finite part of the line for the
variable $s$, and let
\begin{equation}
\tag{7}\label{eq:7}
  f_1(s),f_2(s),f_3(s),\ldots
\end{equation}
be an infinite sequence of functions of the variable $s$ that are continuous
in $J$, whose absolute values remain below a finite bound, and all of whose
difference quotients remain below a finite bound, so that for all $s,t$ in
$J$ and all $h$ one always has
\[
  \left|\frac{f_h(s)-f_h(t)}{s-t}\right|\le E,
\]
where $E$ denotes a finite quantity independent of $h$ and of the choice of
the arguments $s,t$.

Then infinitely many functions
\[
  f_1^*(s),f_2^*(s),f_3^*(s),\ldots
\]
can be selected from this very sequence of functions \eqref{eq:7} in such a
way that the sequence of these functions converges at every point
of $J$ to a finite limiting value, and does so continuously (hence also
uniformly); it follows that the limit
\[
  \mathop{\mathrm{L}}\limits_{p=\infty}f_p^*(s)
\]
represents a function of $s$ continuous in $J$.
\end{lemma}

\begin{proof}
In the manner used in \S5 of my cited paper on Dirichlet's principle, select
from the sequence \eqref{eq:7} a sequence of functions
$f_1^*(s),f_2^*(s),\ldots$ such that, at every point of a countable point set
$P^*$ everywhere dense in $J$, the sequence of these functions converges to a
finite limiting value.  We now show that the functions so selected have the
property required in the lemma.

Indeed, let $s_1,s_2,\ldots$ denote any infinite sequence of points of $J$
converging to a definite point $s$.  If $\varepsilon$ denotes an arbitrarily
small positive quantity, first choose a point $s^*$ of the point set
\originalpage{163}$P^*$ so that
\begin{equation}
\tag{8}\label{eq:8}
  |s^*-s|<\frac{\varepsilon}{8E}.
\end{equation}
Because the sequence of values
\[
  f_1^*(s^*),f_2^*(s^*),\ldots
\]
converges, a number $N$ can then certainly be found such that, for all
indices $p,q$ exceeding $N$,
\begin{equation}
\tag{9}\label{eq:9}
  |f_p^*(s^*)-f_q^*(s^*)|<\frac{\varepsilon}{2}.
\end{equation}
Because the sequence of points $s_1,s_2,\ldots$ converges, this integer $N$
may at the same time be chosen so large that, for all indices $p,q$ exceeding
$N$,
\[
  |s-s_p|<\frac{\varepsilon}{8E},
  \qquad
  |s-s_q|<\frac{\varepsilon}{8E}.
\]
With the aid of \eqref{eq:8}, it follows that, for all indices $p,q$
exceeding $N$,
\begin{equation}
\tag{10}\label{eq:10}
  |s^*-s_p|<\frac{\varepsilon}{4E},
  \qquad
  |s^*-s_q|<\frac{\varepsilon}{4E}.
\end{equation}
On the other hand, by the hypothesis of our lemma,
\[
  \left|\frac{f_p^*(s^*)-f_p^*(s_p)}{s^*-s_p}\right|<E,
  \qquad
  \left|\frac{f_q^*(s^*)-f_q^*(s_q)}{s^*-s_q}\right|<E,
\]
and consequently, with the aid of \eqref{eq:10},
\[
  |f_p^*(s^*)-f_p^*(s_p)|<\frac{\varepsilon}{4},
  \qquad
  |f_q^*(s^*)-f_q^*(s_q)|<\frac{\varepsilon}{4}.
\]
Adding these inequalities to \eqref{eq:9}, we obtain, for all indices $p,q$
exceeding $N$,
\[
\begin{aligned}
  &|f_p^*(s^*)-f_q^*(s^*)|
   +|f_p^*(s^*)-f_p^*(s_p)|
   +|f_q^*(s^*)-f_q^*(s_q)|<\varepsilon,
\end{aligned}
\]
and all the more
\[
  |f_p^*(s_p)-f_q^*(s_q)|<\varepsilon.
\]
Since $\varepsilon$ was arbitrarily small, it follows that the sequence of
values
\[
  f_1^*(s_1),f_2^*(s_2),\ldots
\]
converges to a finite limiting value; that is, the sequence of functions
\[
  f_1^*(s),f_2^*(s),\ldots
\]
\originalpage{164}converges continuously at every point $s$ of the interval $J$.  This proves
the lemma completely.
\end{proof}

We return to the theory of quadratic forms, and in what follows we first
assume that $\lambda=\infty$ is not a condensation value of $K$; that is,
$K$ is to be a form such that all eigenvalues of $K_n$, for every $n$, lie in
the finite interval $J$.  We then define certain functions of the variable
$\lambda$ belonging to $K_n$ as follows: in general let
\begin{equation}
\tag{11}\label{eq:11}
\left.
\begin{aligned}
  \chi_p^{(n)}(\lambda)&=0
    &&\text{for }\lambda\le\lambda_p^{(n)},\\
  \chi_p^{(n)}(\lambda)&=(L_p^{(n)}(x))^2
      (\lambda-\lambda_p^{(n)})
    &&\text{for }\lambda>\lambda_p^{(n)},
\end{aligned}
\right\}
\qquad p=1,\ldots,n,
\end{equation}
and put
\[
  \chi^{(n)}(\lambda)
  =\chi_1^{(n)}(\lambda)+\cdots+\chi_n^{(n)}(\lambda).
\]
Then the $\chi_p^{(n)}$, and consequently also $\chi^{(n)}$, are continuous
functions of $\lambda$ which are nowhere negative and never decrease when
the argument $\lambda$ increases while $n$ is held fixed.  The difference
quotients of these functions, namely
\[
  \frac{\chi_p^{(n)}(\lambda)-\chi_p^{(n)}(\mu)}{\lambda-\mu},
  \qquad
  \frac{\chi^{(n)}(\lambda)-\chi^{(n)}(\mu)}{\lambda-\mu},
\]
are likewise, as we see at once, nowhere negative and never decrease if one
of the arguments $\lambda,\mu$ is held fixed while the other increases.  More
precisely, the equations
\[
  \frac{\chi_p^{(n)}(\lambda)-\chi_p^{(n)}(\mu)}{\lambda-\mu}
  =\pi_p(L_p^{(n)}(x))^2
\]
and consequently
\[
  \frac{\chi^{(n)}(\lambda)-\chi^{(n)}(\mu)}{\lambda-\mu}
  =\pi_1(L_1^{(n)}(x))^2+\cdots+\pi_n(L_n^{(n)}(x))^2
\]
hold, where $\pi_1,\ldots,\pi_n$ denote quantities depending on
$\lambda,\mu$ and lying between $0$ and $1$.  This equation immediately
implies
\[
  \left|
  \frac{\chi^{(n)}(\lambda)-\chi^{(n)}(\mu)}{\lambda-\mu}
  \right|
  \le (L_1^{(n)}(x))^2+\cdots+(L_n^{(n)}(x))^2,
\]
or, by \eqref{eq:5},
\begin{equation}
\tag{12}\label{eq:12}
  \left|
  \frac{\chi^{(n)}(\lambda)-\chi^{(n)}(\mu)}{\lambda-\mu}
  \right|
  \le x_1^2+\cdots+x_n^2.
\end{equation}

\originalpage{165}
Finally, we observe that on every interval containing none of the eigenvalues
$\lambda_p^{(n)}$ of $K_n$, the functions $\chi_p^{(n)}(\lambda)$, and hence
also $\chi^{(n)}(\lambda)$, are linear functions of $\lambda$.  For every
$\lambda$ lying outside $J$ on its negative side,
$\chi^{(n)}(\lambda)$ is identically zero; hence, with regard to
\eqref{eq:12}, one always has
\[
  \chi^{(n)}(\lambda)\le (x_1^2+\cdots+x_n^2)J,
\]
where $J$ denotes the length of the interval $J$.  On the positive side
outside $J$,
\[
  \chi^{(n)}(\lambda)
  =\lambda(x_1^2+\cdots+x_n^2)+C(x),
\]
where $C(x)$ denotes a form independent of $\lambda$.  With respect to
$x_1,\ldots,x_n$, the expression $\chi^{(n)}$ is a quadratic form.

Now let infinitely many variables $x_1,x_2,\ldots$ be given.  We consider
those special systems of values which arise when any one of these variables
is set equal to $1$ and all the others equal to $0$, or when any two of these
variables are set equal to $1/\sqrt2$ and all the others equal to $0$.  We
imagine all these special systems of values of the infinitely many variables
$x_1,x_2,\ldots$ arranged in a definite sequence, and call them, in this
order, the first, second, third, and so on, special system of values of the
infinitely many variables $x_1,x_2,\ldots$.  By \eqref{eq:12}, for all these
special systems of values the inequality
\begin{equation}
\tag{13}\label{eq:13}
  \left|
  \frac{\chi^{(n)}(\lambda)-\chi^{(n)}(\mu)}{\lambda-\mu}
  \right|\le 1
\end{equation}
holds.

We now consider the sequence of forms
\begin{equation}
\tag{14}\label{eq:14}
  \chi^{(1)}(\lambda),\quad
  \chi^{(2)}(\lambda),\quad
  \chi^{(3)}(\lambda),\quad\ldots
\end{equation}
and apply \hyperref[lem:1]{Lemma 1}, with the first special system of values
understood to have been substituted for the variables $x_1,x_2,\ldots$.
We see that, in the interval $J$, the values of all the forms in the sequence
\eqref{eq:14} must lie below a finite bound, and that by \eqref{eq:13} the
remaining hypothesis of our lemma is also fulfilled.  By that lemma, an
infinite sequence of forms
\begin{equation}
\tag{15}\label{eq:15}
  \chi^{(1^*)}(\lambda),\quad
  \chi^{(2^*)}(\lambda),\quad
  \chi^{(3^*)}(\lambda),\quad\ldots
\end{equation}
can therefore certainly be selected from \eqref{eq:14} such that, when the
first special system of values is substituted therein for the variables
$x_1,x_2,\ldots$, it converges uniformly in the interval $J$ to a certain
continuous function of $\lambda$.

\originalpage{166}
We next apply our lemma to the sequence of forms \eqref{eq:15}.  According to
the lemma, there can certainly be selected from that sequence an infinite
sequence
\begin{equation}
\tag{16}\label{eq:16}
  \chi^{(1^{**})}(\lambda),\quad
  \chi^{(2^{**})}(\lambda),\quad
  \chi^{(3^{**})}(\lambda),\quad\ldots
\end{equation}
which, when the second special system of values is substituted therein for
the variables $x_1,x_2,\ldots$, converges uniformly to a certain continuous
function of $\lambda$.  Applying our lemma further to the sequence of
functions \eqref{eq:16}, we arrive by selection at a sequence of functions
\[
  \chi^{(1^{***})}(\lambda),\quad
  \chi^{(2^{***})}(\lambda),\quad
  \chi^{(3^{***})}(\lambda),\quad\ldots,
\]
which, after substitution of the third special system of values, converges
uniformly to a certain continuous function of $\lambda$.

Finally, imagining the procedure continued without limit, we consider the
sequence of forms
\begin{equation}
\tag{17}\label{eq:17}
  \chi^{(1^*)}(\lambda),\quad
  \chi^{(2^{**})}(\lambda),\quad
  \chi^{(3^{***})}(\lambda),\quad\ldots;
\end{equation}
it converges uniformly in $\lambda$ when the first, the second, the third,
and so on, special system of values is substituted therein for the variables
$x_1,x_2,\ldots$.  But the coefficient of $x_p x_q$ in a quadratic form can
in general be expressed as a linear combination of the values that the
quadratic form assumes for
\[
  x_p,x_q=1,0;\qquad 0,1;\qquad
  \frac1{\sqrt2},\frac1{\sqrt2},
\]
respectively, all other variables being set equal to zero.  It follows that,
in general, the coefficient of $x_p x_q$ in the sequence of forms
\eqref{eq:17} also converges to a certain function
$\chi_{pq}(\lambda)$ continuous in $\lambda$.  We put
\[
  \chi(\lambda)
  =\sum_{p,q=1,2,\ldots}\chi_{pq}(\lambda)x_p x_q,
\]
so that $\chi(\lambda)$ denotes a quadratic form in the infinitely many
variables $x_1,x_2,\ldots$, whose coefficients are continuous functions of
$\lambda$.  These functions of $\lambda$ are initially defined only within
$J$; but since, in place of $J$, we may choose an arbitrarily large interval
containing $J$, the definition of these functions is thereby given for all
finite values of $\lambda$.

The values of any truncation of a quadratic form in infinitely many variables
can be represented as linear combinations of the values which the quadratic
form assumes for our special systems of values.  It follows that if,
\originalpage{167}in the forms of the sequence \eqref{eq:17}, we leave
$x_1,\ldots,x_n$ arbitrary and set all the remaining variables
$x_{n+1},x_{n+2},\ldots$ equal to zero---that is, take the same truncation of
every form in \eqref{eq:17}---the resulting sequence of truncations likewise
converges uniformly, and indeed to the value assumed by the corresponding
truncation of $\chi(\lambda)$.  If we also put
\[
  1^*=m_1,\qquad 2^{**}=m_2,\qquad 3^{***}=m_3,\quad\ldots,
\]
then the equation
\[
  \mathop{\mathrm{L}}\limits_{h=\infty}
  [\chi^{(m_h)}(\lambda)]=[\chi(\lambda)]
\]
holds.

For brevity and clarity, in what follows we shall omit square brackets in an
equation or inequality.  Thus, according to this convention, an equation or
inequality between forms in infinitely many variables is always to be
understood as holding identically for all the variables when the same
truncations of the forms are taken on both sides of the formula.  The last
equation then reads
\begin{equation}
\tag{18}\label{eq:18}
  \mathop{\mathrm{L}}\limits_{h=\infty}
  \chi^{(m_h)}(\lambda)=\chi(\lambda).
\end{equation}

At the same time we observe that the properties named above of the functions
$\chi^{(n)}(\lambda)$, regarded as functions of $\lambda$, carry over at once
to an arbitrary truncation $[\chi(\lambda)]$ of the form $\chi(\lambda)$ when
we regard this truncation as a function of $\lambda$.  We thus see that the
function $[\chi(\lambda)]$ is likewise nowhere negative and never decreases
as the argument $\lambda$ increases.  Further, the difference quotient of
this function, that is, the expression
\[
  \frac{[\chi(\lambda)]-[\chi(\mu)]}{\lambda-\mu},
\]
is again nowhere negative and, if one of the arguments $\lambda,\mu$ remains
fixed while the other increases, never decreases.  By \eqref{eq:12}, the
inequality
\begin{equation}
\tag{19}\label{eq:19}
  \frac{\chi(\lambda)-\chi(\mu)}{\lambda-\mu}\le(x,x)
\end{equation}
holds for this difference quotient.  Within an interval containing no
condensation value of the form $K$, $[\chi(\lambda)]$ becomes a linear
function of $\lambda$.  Outside $J$,\originalpage{168} on the negative side,
$[\chi(\lambda)]$ is identically zero, and on the positive side it equals
$\lambda[(x,x)]+C(x)$.

We now imagine the first special system of values substituted in the form
$\chi(\lambda)$, and denote the resulting function of $\lambda$ by
$\chi(\lambda)_1$.  According to the foregoing discussion, if $\lambda$ is
held fixed and $\mu$ is assigned a value exceeding $\lambda$, the difference
quotient of $\chi(\lambda)_1$ certainly will not increase as $\mu$ decreases
toward $\lambda$, and hence, as $\mu$ falls to $\lambda$, will tend to a
limiting value.  Thus $\chi(\lambda)_1$ certainly possesses a forward
derivative at every $\lambda$; denote it by
$\mathfrak f^{(+)}(\lambda)_1$.  In the same way it is shown that
$\chi(\lambda)_1$ possesses a backward derivative at every $\lambda$; denote
it by $\mathfrak f^{(-)}(\lambda)_1$.

From the facts stated above concerning the difference quotients of
$\chi(\lambda)_1$, it follows at the same time that both
$\mathfrak f^{(+)}(\lambda)_1$ and $\mathfrak f^{(-)}(\lambda)_1$ are
functions of $\lambda$ which are nowhere negative and do not decrease as the
argument $\lambda$ increases, and for which, moreover, one always has
\begin{equation}
\tag{20}\label{eq:20}
\begin{aligned}
  \mathfrak f^{(-)}(\lambda)_1
    &\le\mathfrak f^{(+)}(\mu)_1 &&\text{for }\lambda\le\mu,\\
  \mathfrak f^{(+)}(\lambda)_1
    &\le\mathfrak f^{(-)}(\mu)_1 &&\text{for }\lambda<\mu.
\end{aligned}
\end{equation}

Furthermore, since by \eqref{eq:13} the derivatives
$\mathfrak f^{(+)}(\lambda)_1$ and $\mathfrak f^{(-)}(\lambda)_1$ cannot
exceed the value $1$, the same is true all the more of the difference between
the forward and backward derivatives at the same point $\lambda$.  Hence, if
$m$ denotes any integer, there are at most $m$ points $\lambda$ for which
\[
  \mathfrak f^{(+)}(\lambda)_1-
  \mathfrak f^{(-)}(\lambda)_1\ge\frac1m.
\]
It follows that the set of values $\lambda$ for which the forward and
backward derivatives differ is necessarily countable.

We now imagine the second special system of values substituted in
$\chi(\lambda)$, proceed with the resulting function $\chi(\lambda)_2$ as we
did above with $\chi(\lambda)_1$, and seek the set of points $\lambda$ at
which the forward and backward derivatives differ.  This set of points is
again certainly countable.  By using the third special system of values we
obtain correspondingly a countable set of points $\lambda$, and so on.  The
totality of all such
\originalpage{169}points is again countable; let them be called the
\emph{eigenvalues} of the form $K$ and be denoted by
$\lambda_1,\lambda_2,\ldots$.  The points $\lambda_1,\lambda_2,\ldots$ and
their condensation points are certainly condensation values of the form
$K$, since $\chi(\lambda)_1,\chi(\lambda)_2,\ldots$ are linear functions of
$\lambda$ outside the condensation values, and consequently the forward and
backward derivatives there are equal.  The totality of the points
$\lambda_1,\lambda_2,\ldots$ shall be called the \emph{point spectrum} or the
\emph{discontinuous spectrum} of the form $K$.

Since the coefficients of a quadratic form can be composed linearly from the
values which the quadratic form assumes on the sequence of special systems of
values of the variables, we conclude that all the coefficients
$\chi_{pq}$ of the form $\chi(\lambda)$ likewise possess both forward and
backward derivatives, and that these are equal at every point except the
points $\lambda_1,\lambda_2,\ldots$.  Let the forward and backward derivatives
of $\chi_{pq}$ be denoted by $\mathfrak f_{pq}^{(+)}$ and
$\mathfrak f_{pq}^{(-)}$, respectively.  At every point $\lambda$ distinct
from $\lambda_1,\lambda_2,\ldots$, these two derivatives agree; there we put
\[
  \mathfrak f_{pq}=\mathfrak f_{pq}^{(+)}=\mathfrak f_{pq}^{(-)}.
\]
Let the quadratic forms with coefficients $\mathfrak f_{pq}^{(+)}$ and
$\mathfrak f_{pq}^{(-)}$ be denoted by $\mathfrak f^{(+)}(\lambda)$ and
$\mathfrak f^{(-)}(\lambda)$, respectively.  At every point distinct from
$\lambda_1,\lambda_2,\ldots$, we put
\[
  \mathfrak f(\lambda)
  =\mathfrak f^{(+)}(\lambda)=\mathfrak f^{(-)}(\lambda).
\]

For the point $\lambda_h$, form the difference
$\mathfrak f_{pq}^{(+)}-\mathfrak f_{pq}^{(-)}$ and take this difference as
the coefficient of $x_p x_q$.  The resulting quadratic form in infinitely
many variables, whose coefficients certainly do not all vanish, shall be
called the \emph{quadratic eigenform of $K$ belonging to the eigenvalue
$\lambda_h$}, and shall be denoted by $E_h$.  Evidently, for every eigenvalue
$\lambda_p$,
\begin{equation}
\tag{21}\label{eq:21}
  \mathfrak f^{(+)}(\lambda_p)
  -\mathfrak f^{(-)}(\lambda_p)=E_p,
\end{equation}
and since, analogously to \eqref{eq:20}, one must also have
\begin{equation}
\tag{22}\label{eq:22}
\begin{aligned}
  \mathfrak f^{(-)}(\lambda)&\le\mathfrak f^{(+)}(\mu)
    &&\text{for }\lambda\le\mu,\\
  \mathfrak f^{(+)}(\lambda)&\le\mathfrak f^{(-)}(\mu)
    &&\text{for }\lambda<\mu,
\end{aligned}
\end{equation}
it certainly follows that
\begin{equation}
\tag{23}\label{eq:23}
  E_p\ge0.
\end{equation}

\originalpage{170}
On the other hand, it follows from \eqref{eq:19}, if we choose for $\mu$ a
value exceeding $\lambda$ and let it converge to $\lambda$, that
\begin{equation}
\tag{24}\label{eq:24}
  \mathfrak f^{(+)}(\lambda)\le(x,x).
\end{equation}

We now consider any $m$ eigenvalues of the form $K$, arranged in increasing
order, say
\[
  \lambda_{p_1},\ldots,\lambda_{p_m},
\]
and the eigenforms belonging to them,
\[
  E_{p_1},\ldots,E_{p_m}.
\]
By \eqref{eq:21} we have
\begin{equation}
\tag{25}\label{eq:25}
  \mathfrak f^{(+)}(\lambda_{p_h})
  -\mathfrak f^{(-)}(\lambda_{p_h})=E_{p_h},
  \qquad(h=1,2,\ldots,m),
\end{equation}
and by \eqref{eq:22},
\begin{equation}
\tag{26}\label{eq:26}
\begin{aligned}
  \mathfrak f^{(-)}(\lambda_{p_2})
    -\mathfrak f^{(+)}(\lambda_{p_1})&\ge0,\\
  &\ \vdots\\
  \mathfrak f^{(-)}(\lambda_{p_m})
    -\mathfrak f^{(+)}(\lambda_{p_{m-1}})&\ge0.
\end{aligned}
\end{equation}
Adding \eqref{eq:25} and \eqref{eq:26} gives
\begin{equation}
\tag{27}\label{eq:27}
  \mathfrak f^{(+)}(\lambda_{p_m})
  -\mathfrak f^{(-)}(\lambda_{p_1})
  \ge E_{p_1}+\cdots+E_{p_m},
\end{equation}
and, since $[\mathfrak f^{(-)}(\lambda)]$ is nowhere negative, this inequality,
together with \eqref{eq:24}, yields the further inequality
\begin{equation}
\tag{28}\label{eq:28}
  E_{p_1}+\cdots+E_{p_m}\le(x,x).
\end{equation}

From \eqref{eq:23} and \eqref{eq:28} we see that the sum
$\sum_{(p)}[E_p]$, extended over all indices $p$, converges, and indeed to a
quadratic form in the $n$ variables $x_1,\ldots,x_n$ that is at most
$x_1^2+\cdots+x_n^2$; that is,
\[
  \sum_{(p)}E_p\le(x,x).
\]
We now define the following forms in the infinitely many variables
$x_1,x_2,\ldots$:
\[
  e(\lambda)=\sum_{(\lambda_p<\lambda)}E_p,
  \qquad
  \eta(\lambda)=\sum_{(\lambda_p<\lambda)}E_p(\lambda-\lambda_p),
\]
where the sums on the right are to be extended over all indices $p$ for which
$\lambda_p<\lambda$.  The coefficients of
\originalpage{171}$\eta(\lambda)$ are continuous functions of $\lambda$.
The truncations $[e(\lambda)]$ and $[\eta(\lambda)]$ of these forms are
functions of $\lambda$ which, as follows without difficulty from
\eqref{eq:23} and the convergence of $\sum_{(p)}[E_p]$, possess the following
properties.

$[e(\lambda)]$ is nowhere negative and never decreases as $\lambda$
increases; outside $J$ it is zero on the negative side, and on the positive
side it is equal to $\sum_{(p)}[E_p]$.

$[e(\lambda)]$ is continuous at every point that is not an eigenvalue; at the
eigenvalue $\lambda_p$ it has a finite jump, and for positive $\tau$
decreasing to zero one has
\[
  \mathop{\mathrm{L}}\limits_{\tau=0}e(\lambda_p-\tau)=e(\lambda_p),
  \qquad
  \mathop{\mathrm{L}}\limits_{\tau=0}e(\lambda_p+\tau)
  =e(\lambda_p)+E_p.
\]

$[\eta(\lambda)]$ has both a forward and a backward derivative; these agree
at every point that is not an eigenvalue and have the value $[e(\lambda)]$.
At the eigenvalue $\lambda_p$ the backward derivative is $[e(\lambda_p)]$ and
the forward derivative is $[e(\lambda_p)]+[E_p]$, so that the excess of the
forward derivative over the backward derivative amounts to $[E_p]$.

From \eqref{eq:27} it follows, if $\lambda'$ and $\lambda''$, with
$\lambda''>\lambda'$, are not eigenvalues, that
\[
  \mathfrak f(\lambda'')-\mathfrak f(\lambda')
  \ge\sum_{(\lambda'<\lambda_p<\lambda'')}E_p;
\]
but
\[
  \sum_{(\lambda'<\lambda_p<\lambda'')}E_p
  =e(\lambda'')-e(\lambda'),
\]
and consequently also
\begin{equation}
\tag{29}\label{eq:29}
  \mathfrak f(\lambda'')-\mathfrak f(\lambda')
  \ge e(\lambda'')-e(\lambda').
\end{equation}
If one of the quantities $\lambda',\lambda''$ is an eigenvalue, or if both
are eigenvalues of $K$, a corresponding inequality follows in the same way
in place of \eqref{eq:29}.  Now put
\[
  \varrho(\lambda)=\chi(\lambda)-\eta(\lambda).
\]
Then, with regard to \eqref{eq:21}, the function $[\varrho(\lambda)]$ has
forward and backward derivatives which agree at every point $\lambda$; this
derivative is moreover, as follows from \eqref{eq:29} or the corresponding
inequality, a function of $\lambda$ that does not decrease as the argument
increases.  This derivative therefore represents a function which is
continuous in $\lambda$.  \originalpage{172}Thus, if we put
\begin{equation}
\tag{30}\label{eq:30}
  \sigma(\lambda)=\sum\sigma_{pq}x_p x_q
  =\frac{d\varrho(\lambda)}{d\lambda}
  =\mathfrak f(\lambda)-e(\lambda),
\end{equation}
then every truncation of the form $\sigma(\lambda)$ is a continuous,
nonnegative function of $\lambda$ that does not decrease as $\lambda$
increases.  The form $\sigma(\lambda)$ in the infinitely many variables,
whose coefficients $\sigma_{pq}$ are continuous functions of $\lambda$, shall
be called the \emph{spectral form} of $K$.

By \eqref{eq:24}, the inequality
\begin{equation}
\tag{31}\label{eq:31}
  \sigma(\lambda)+e(\lambda)\le(x,x)
\end{equation}
holds.

For all values $\lambda$ lying outside $J$ on its negative side, the forms
$\mathfrak f(\lambda),e(\lambda),\sigma(\lambda)$ become identically zero; on
the positive side they pass into certain forms independent of $\lambda$,
which we shall denote respectively by
$\mathfrak f(+\infty),e(+\infty),\sigma(+\infty)$.  By \eqref{eq:30} we have
\begin{equation}
\tag{32}\label{eq:32}
  \sigma(+\infty)=\mathfrak f(+\infty)-e(+\infty).
\end{equation}
Since $[\chi(\lambda)]$ passes, to the right of $J$, into
$\lambda[(x,x)]+C(x)$, we have
\begin{equation}
\tag{33}\label{eq:33}
  \mathfrak f(+\infty)=(x,x),
\end{equation}
and, as observed earlier,
\[
  e(+\infty)=\sum_p E_p.
\]

We now select those real values $\lambda$ in every neighborhood of which
there still exist points $\lambda'$ for which
\[
  \sigma(\lambda)=\sigma(\lambda')
\]
does not hold identically in all the variables $x_1,x_2,\ldots$. The set $s$
of all such points $\lambda$ is perfect (closed and dense in itself); it shall
be called the \emph{interval spectrum} or the \emph{continuous spectrum} of
the form $K$.  Outside the condensation values of $K$, all coefficients of
$\chi(\lambda)$ are linear in $\lambda$, those of $e(\lambda)$ are constant,
and consequently the coefficients of $\sigma(\lambda)$ likewise become
constant; thus the interval spectrum certainly lies within the condensation
values of the form $K$.  The point spectrum together with the accumulation
points of the eigenvalues, and the interval spectrum, taken together, shall
be called the \emph{spectrum} of the form $K$.

If we now put
\[
  \sigma(+\infty)
  =\int_{-\infty}^{+\infty}d\sigma(\lambda)
  =\int_{(s)}d\sigma(\lambda),
\]
\originalpage{173}then from \eqref{eq:32} and \eqref{eq:33} we obtain
\[
  (x,x)=e(+\infty)+\int_{(s)}d\sigma(\lambda),
\]
or
\begin{equation}
\tag{34}\label{eq:34}
  (x,x)=\sum_{(p)}E_p+\int_{(s)}d\sigma(\lambda),
\end{equation}
where the sum is to be extended over all $p$.

We now return to the quadratic form $K_n$ with the finite number $n$ of
variables.  If $\lambda$ denotes a complex quantity, or a real quantity
different from every eigenvalue $\lambda_p^{(n)}$ of the form $K_n$, then
from the definition \eqref{eq:11} of the function $\chi_p^{(n)}$ we obtain at
once
\[
\begin{aligned}
  \int_{-\infty}^{+\infty}
    \frac{\chi_p^{(n)}(\mu)}{(\mu-\lambda)^3}\,d\mu
  &=(L_p^{(n)}(x))^2
    \int_{\lambda_p^{(n)}}^{+\infty}
    \frac{\mu-\lambda_p^{(n)}}{(\mu-\lambda)^3}\,d\mu\\
  &=\frac12\frac{(L_p^{(n)}(x))^2}{\lambda_p^{(n)}-\lambda}.
\end{aligned}
\]
Here the integration is to be carried out along the real axis from
$\mu=-\infty$ to $\mu=+\infty$; when $\lambda$ is real, however, the point
$\lambda$ is to be bypassed by a short detour in the complex $\mu$-plane,
with the observation that, in a neighborhood of the point $\lambda$,
$\chi_p^{(n)}(\mu)$, being a linear function, is defined also for complex
$\mu$.  Summing over $p=1,\ldots,n$, we find
\begin{equation}
\tag{35}\label{eq:35}
  \int_{-\infty}^{+\infty}
    \frac{\chi^{(n)}(\mu)}{(\mu-\lambda)^3}\,d\mu
  =\frac12\left\{
    \frac{(L_1^{(n)}(x))^2}{\lambda_1^{(n)}-\lambda}
    +\cdots+
    \frac{(L_n^{(n)}(x))^2}{\lambda_n^{(n)}-\lambda}
  \right\},
\end{equation}
where the integration is to be carried out as above.  On the other hand,
from \eqref{eq:3} and \eqref{eq:5} it follows that
\begin{equation}
\tag{36}\label{eq:36}
  \frac{K_n(\lambda,x)-(x,x)_n}{\lambda}
  =\frac{(L_1^{(n)}(x))^2}{\lambda_1^{(n)}-\lambda}
   +\cdots+
   \frac{(L_n^{(n)}(x))^2}{\lambda_n^{(n)}-\lambda}.
\end{equation}

If in \eqref{eq:35} and \eqref{eq:36} we put for $n$ the special numbers
$m_1,m_2,\ldots$ and take a fixed truncation on both sides, and if further
$\lambda$ is understood to be any complex quantity or a real quantity that
does not belong to the spectrum of the form $K$, then by \eqref{eq:18}, on
passing to the limit, we obtain the formula
\begin{equation}
\tag{37}\label{eq:37}
  \mathop{\mathrm{L}}\limits_{h=\infty}
  \frac{[K_{m_h}(\lambda,x)]-[(x,x)_{m_h}]}{\lambda}
  =2\int_{-\infty}^{+\infty}
    \frac{[\chi(\mu)]}{(\mu-\lambda)^3}\,d\mu.
\end{equation}
Here again the integration is along the real axis from $\mu=-\infty$ to
$\mu=+\infty$; when $\lambda$ is real, however, the point $\lambda$ is to be
bypassed by a short detour in the complex $\mu$-plane, with the observation
that, in a neighborhood of the point $\lambda$, $\chi(\mu)$, being a linear
function, is defined also\originalpage{174} for complex $\mu$.  Consequently,
the integral on the right represents a function of $\lambda$ which has a
regular analytic character with respect to $\lambda$ for all complex
$\lambda$ and for all real $\lambda$ not belonging to the spectrum of $K$.

We put
\begin{equation}
\tag{38}\label{eq:38}
  K(\lambda,x)=(x,x)+2\lambda
  \int_{-\infty}^{+\infty}
    \frac{\chi(\mu)}{(\mu-\lambda)^3}\,d\mu
\end{equation}
and call this expression the \emph{resolvent of the form $K$}.  Every
coefficient or truncation of it is likewise regular analytic for all complex
$\lambda$ and for the finite real values of $\lambda$ not belonging to the
spectrum of $K$; one also sees that it is regular analytic at
$\lambda=\infty$.

The equation \eqref{eq:37} found above becomes
\[
  \mathop{\mathrm{L}}\limits_{h=\infty}K_{m_h}(\lambda,x)=K(\lambda,x).
\]
From the definition of the form $\eta(\lambda)$ we obtain the equation
\begin{equation}
\tag{39}\label{eq:39}
\begin{aligned}
  \int_{-\infty}^{+\infty}
    \frac{\eta(\mu)}{(\mu-\lambda)^3}\,d\mu
  &=\sum_{(p)}E_p
    \int_{\lambda_p}^{+\infty}
    \frac{\mu-\lambda_p}{(\mu-\lambda)^3}\,d\mu\\
  &=\frac12\sum_{(p)}\frac{E_p}{\lambda_p-\lambda},
\end{aligned}
\end{equation}
where $\lambda$ is to be understood as any complex quantity, or as a real
quantity not belonging to the spectrum of the form $K$; when $\lambda$ is
real, the integration with respect to $\mu$ is to be carried out with a short
detour around the point $\lambda$ in the complex $\mu$-plane.

Since
\[
  \varrho(\lambda)=\chi(\lambda)-\eta(\lambda),
  \qquad
  \sigma(\lambda)=\frac{d\varrho(\lambda)}{d\lambda},
\]
we obtain
\[
\begin{aligned}
  &\int_{-\infty}^{+\infty}
    \frac{\chi(\mu)}{(\mu-\lambda)^3}\,d\mu
  -\int_{-\infty}^{+\infty}
    \frac{\eta(\mu)}{(\mu-\lambda)^3}\,d\mu\\
  &\qquad=\int_{-\infty}^{+\infty}
    \frac{\varrho(\mu)}{(\mu-\lambda)^3}\,d\mu
  =\frac12\int_{-\infty}^{+\infty}
    \frac{\sigma(\mu)}{(\mu-\lambda)^2}\,d\mu
  =\frac12\int_{(s)}\frac{d\sigma(\mu)}{\mu-\lambda}.
\end{aligned}
\]
From \eqref{eq:38} and \eqref{eq:39} it follows that
\begin{equation}
\tag{40}\label{eq:40}
  K(\lambda,x)
  =(x,x)+\lambda\sum_{(p)}\frac{E_p}{\lambda_p-\lambda}
  +\lambda\int_{(s)}\frac{d\sigma(\mu)}{\mu-\lambda},
\end{equation}
\originalpage{175}and hence, with regard to \eqref{eq:34},
\[
  K(\lambda,x)
  =\sum_{(p)}\frac{E_p}{1-\lambda/\lambda_p}
  +\int_{(s)}\frac{d\sigma(\mu)}{1-\lambda/\mu}.
\]

This formula is the sought analogue of the partial-fraction representation
\eqref{eq:3} of the resolvent $K_n(\lambda,x)$ of the form $K_n$ with the
finite number $n$ of variables.

We summarize the results obtained as follows.

\begin{theorem}[I]\label{thm:I}
The resolvent of a quadratic form $K$ for which $\lambda=\infty$ is not a
condensation value is a quadratic form in infinitely many variables,
\[
  K(\lambda,x)=\sum_{p,q}K_{pq}x_p x_q,
\]
whose coefficients are regular analytic functions of the argument $\lambda$
for all arguments $\lambda$ lying outside the spectrum of the form $K$.

If $m_1,m_2,\ldots$ is a certain sequence of integers increasing to infinity,
then for every truncation of the resolvent the equation
\begin{equation}
\tag{41}\label{eq:41}
  \mathop{\mathrm{L}}\limits_{h=\infty}K_{m_h}(\lambda,x)=K(\lambda,x)
\end{equation}
holds, where $K_{m_h}$ denotes the resolvent of the form $K_{m_h}$.

The resolvent $K$ admits the following representation (for every truncation):
\begin{equation}
\tag{42}\label{eq:42}
  K(\lambda,x)
  =\sum_{(p)}\frac{E_p}{1-\lambda/\lambda_p}
  +\int_{(s)}\frac{d\sigma(\mu)}{1-\lambda/\mu}.
\end{equation}
Here the sum is to be extended over the entire point spectrum of $K$, that
is, over all eigenvalues of $K$; $E_p$ denotes in general the quadratic
eigenform belonging to $\lambda_p$.  It is a form whose truncations are not
negative for any values of the variables $x_1,x_2,\ldots$.  The integral is
to be extended over the interval spectrum of $K$.  The spectral form
$\sigma(\lambda)$ is a form whose coefficients are continuous functions of
$\lambda$, and whose truncations, as the argument $\lambda$ increases within
the interval spectrum $s$, are nondecreasing functions of $\lambda$ that are
not all constant, but in every interval lying outside $s$ are all constant.
Finally, the equation (for every truncation)
\begin{equation}
\tag{43}\label{eq:43}
  (x,x)=\sum_{(p)}E_p+\int_{(s)}d\sigma(\lambda)
\end{equation}
holds.
\end{theorem}

\originalpage{176}
If the absolute values of all truncations of a quadratic or bilinear form in
infinitely many variables $x_1,x_2,\ldots,y_1,y_2,\ldots$ lie below a finite
bound independent of the choice of truncation whenever the variables are
subjected to the conditions
\begin{equation}
\tag{44}\label{eq:44}
  (x,x)\le1,\qquad(y,y)\le1,
\end{equation}
then the quadratic or bilinear form, respectively, shall be called a
\emph{bounded form}.

The bilinear form belonging to a bounded quadratic form is likewise a bounded
form.

For example, by \eqref{eq:23} and \eqref{eq:28} the eigenforms $E_p$ are
always bounded forms; and since by \eqref{eq:31} the inequality
\[
  \sigma(\lambda)+e(\lambda)\le(x,x)
\]
holds, the spectral form $\sigma(\lambda)$, as well as the forms
$\mathfrak f(\lambda)$ and $e(\lambda)$, are bounded forms.  Finally, if
$\lambda$ denotes a value lying outside the spectrum of $K(x)$, the resolvent
$K(\lambda,x)$ is a bounded form; indeed, the sum component and the integral
component of $K(\lambda,x)$ are separately bounded forms.

The concepts ``truncation'' and ``bounded'' can evidently be applied in the
same sense also to linear forms in infinitely many variables.  A linear form
\[
  L(x)=l_1x_1+l_2x_2+\cdots
\]
is a bounded form if and only if the sum of the squares of its coefficients,
\[
  l_1^2+l_2^2+\cdots,
\]
is finite.

The truth of this is easily seen with the aid of the following two facts.

\textup{I.} If $u_1,u_2,\ldots$ and $v_1,v_2,\ldots$ are arbitrary
quantities, then the inequality
\[
  (u,v)^2\le(u,u)(v,v)
\]
always holds.

\textup{II.} If $u_1,u_2,\ldots$ are arbitrary quantities, $M$ is a finite
positive number, and moreover the inequality
\[
  (u,x)\le M
\]
holds for all values $x_1,x_2,\ldots$ satisfying condition \eqref{eq:44},
then always
\[
  (u,u)\le M^2.
\]

\originalpage{177}
Henceforth we shall throughout consider only such systems of values of the
infinitely many variables $x_1,x_2,\ldots,y_1,y_2,\ldots$, and so on, as
satisfy condition \eqref{eq:44}.  If $a_1,a_2,\ldots$ is such a system of
values, then, when $L(x)$ is a bounded linear form, the limit of its $n$th
truncation $[L(a)]_n$ certainly exists as $n\to\infty$; we put
\[
  L(a)=\mathop{\mathrm{L}}\limits_{n=\infty}[L(a)]_n.
\]
We have thus seen that a bounded linear form in the infinitely many variables
$x_1,x_2,\ldots$ assumes a definite finite value for every system of values
of these variables under consideration, and consequently represents a
\emph{function} of the infinitely many variables $x_1,x_2,\ldots$.

In general, a function $F(x_1,x_2,\ldots)$ of the infinitely many variables
$x_1,x_2,\ldots$ shall be called \emph{continuous} at the point
$a_1,a_2,\ldots$ if the value of
$F(a_1+\varepsilon_1,a_2+\varepsilon_2,\ldots)$ converges to the value of
$F(a_1,a_2,\ldots)$ as the sum of the squares of
$\varepsilon_1,\varepsilon_2,\ldots$ decreases to zero; that is, if
\[
  \mathop{\mathrm{L}}\limits_{(\varepsilon_1^2+\varepsilon_2^2+\cdots=0)}
  F(a_1+\varepsilon_1,a_2+\varepsilon_2,\ldots)
  =F(a_1,a_2,\ldots).
\]
We see at once that the bounded linear form $L(x)$ represents a continuous
function of the infinitely many variables.

The corresponding statements also hold for a bounded bilinear form.  To see
this, let $A(x,y)$ denote a bounded bilinear form and let
$a_1,a_2,\ldots$ and $b_1,b_2,\ldots$ be particular systems of values of the
infinitely many variables.  We then put
\begin{align*}
  x_1&=a_1,\ldots,x_n=a_n,
    \quad x_{n+1}=0,\quad x_{n+2}=0,\ldots,\\
  x'_1&=0,\ldots,x'_n=0,\\
  x'_{n+1}&=\frac{a_{n+1}}{\alpha_{nm}},
    \quad x'_{n+2}=\frac{a_{n+2}}{\alpha_{nm}},\ldots,
    \quad x'_m=\frac{a_m}{\alpha_{nm}},\\
  x'_{m+1}&=0,\quad x'_{m+2}=0,\ldots,\\
  y_1&=b_1,\ldots,y_n=b_n,
    \quad y_{n+1}=0,\quad y_{n+2}=0,\ldots,\\
  y'_1&=0,\ldots,y'_n=0,\\
  y'_{n+1}&=\frac{b_{n+1}}{\beta_{nm}},
    \quad y'_{n+2}=\frac{b_{n+2}}{\beta_{nm}},\ldots,
    \quad y'_m=\frac{b_m}{\beta_{nm}},\\
  y'_{m+1}&=0,\quad y'_{m+2}=0,\ldots,
\end{align*}
where
\[
  \alpha_{nm}=\sqrt{a_{n+1}^2+a_{n+2}^2+\cdots+a_m^2},
  \qquad
  \beta_{nm}=\sqrt{b_{n+1}^2+b_{n+2}^2+\cdots+b_m^2},
\]
and where $n$ and $m>n$ denote arbitrary integers.  If $[A]_m$ and $[A]_n$
denote the corresponding truncations of $A$, then
\originalpage{178}
\[
\begin{aligned}
  \relax[A(a,b)]_m
    &=A(x+\alpha_{nm}x',y+\beta_{nm}y')\\
    &=A(x,y)+\alpha_{nm}A(x',y)+\beta_{nm}A(x,y')
      +\alpha_{nm}\beta_{nm}A(x',y')\\
    &=[A(a,b)]_n+\alpha_{nm}A(x',y)+\beta_{nm}A(x,y')
      +\alpha_{nm}\beta_{nm}A(x',y').
\end{aligned}
\]
Since the absolute values of $A(x',y),A(x,y'),A(x',y')$ remain below a bound
independent of $n,m$, while moreover $\alpha_{nm},\beta_{nm}$ vanish as
$n,m$ increase, it follows that $[A(a,b)]_n$ must converge to a limiting
value as $n$ increases; denote this value briefly by $A(a,b)$.  In an
analogous manner it follows that the function of the infinitely many
variables $x_1,x_2,\ldots$ and $y_1,y_2,\ldots$ represented by $A(x,y)$ is
continuous.

In particular, we conclude that a bounded quadratic form in infinitely many
variables $x_1,x_2,\ldots$ always represents a continuous function of them.

From the statements just proved we further infer the following fact.  If on
both sides of a formula there stand a finite number of bounded quadratic or
bilinear forms in infinitely many variables, and the equation or inequality
represented by the formula is valid for all truncations of both sides, then
it is valid in general for all systems of values of the infinitely many
variables---where only those systems of values satisfying conditions
\eqref{eq:44} are always meant.

If, in the bounded bilinear form $A(x,y)$, all the variables
$y_1,y_2,\ldots$ are set equal to zero with the sole exception of the one
variable $y_p$, which is assigned the value $1$, there results a bounded
linear form in the variables $x_1,x_2,\ldots$; denote it by
\[
  \frac{\partial A(x,y)}{\partial y_p}.
\]
Since $A(x,y)$ is bounded, it also follows that
\[
  y_1\frac{\partial[A(x,y)]_n}{\partial y_1}
  +\cdots+
  y_m\frac{\partial[A(x,y)]_n}{\partial y_m}
\]
lies in absolute value below a finite bound independent of $n$ and $m$; and,
with regard to fact II above, therefore also
\[
  \left(\frac{\partial[A(x,y)]_n}{\partial y_1}\right)^2
  +\cdots+
  \left(\frac{\partial[A(x,y)]_n}{\partial y_m}\right)^2
\]
lies below a bound independent of $n$ and $m$.  Taking first $n\to\infty$ and
then $m\to\infty$, we see that the
\originalpage{179}sum of squares
\[
  \left(\frac{\partial A(x,y)}{\partial y_1}\right)^2
  +\left(\frac{\partial A(x,y)}{\partial y_2}\right)^2+\cdots
\]
converges to a finite quantity.  If $B(x,y)$ is likewise a bounded bilinear
form, then the sum of squares
\[
  \left(\frac{\partial B(x,y)}{\partial x_1}\right)^2
  +\left(\frac{\partial B(x,y)}{\partial x_2}\right)^2+\cdots
\]
also remains below a finite bound; and consequently, with regard to fact I
above, the infinite series
\[
  \frac{\partial A(x,y)}{\partial y_1}
  \frac{\partial B(x,y)}{\partial x_1}
  +\frac{\partial A(x,y)}{\partial y_2}
  \frac{\partial B(x,y)}{\partial x_2}+\cdots
\]
must converge absolutely.  It then necessarily represents once more a
bilinear form in the variables
$x_1,x_2,\ldots,y_1,y_2,\ldots$, which we denote by
\[
  A(x,\mathord{\cdot})B(\mathord{\cdot},y)
\]
and call the \emph{convolution} of the bilinear forms $A,B$.

The convolution is certainly a bounded bilinear form; with regard to facts I
and II above, the preceding discussion establishes the following lemma.

\begin{lemma}[2]\label{lem:2}
If $M,N$ denote two positive constants such that, for all
$x_1,x_2,\ldots$ and $y_1,y_2,\ldots$,
\[
  |A(x,y)|\le M,
  \qquad
  |B(x,y)|\le N,
\]
then the convolution necessarily satisfies the inequality
\[
  |A(x,\mathord{\cdot})B(\mathord{\cdot},y)|\le MN.
\]
\end{lemma}

At the same time we state the following lemmas, whose truth is immediately
evident.

\begin{lemma}[3]\label{lem:3}
For every truncation of the convolution of two bounded bilinear forms $A,B$,
\[
  [A(x,\mathord{\cdot})B(\mathord{\cdot},y)]_m
  =\mathop{\mathrm{L}}\limits_{n=\infty}
   [A_n(x,\mathord{\cdot})B_n(\mathord{\cdot},y)]_m,
\]
where on the right the expression under the limit is the $m$th truncation of
the convolution of the $n$th truncations of $A,B$.  The value of the
convolution is consequently
\[
  A(x,\mathord{\cdot})B(\mathord{\cdot},y)
  =\mathop{\mathrm{L}}\limits_{m=\infty}
   \mathop{\mathrm{L}}\limits_{n=\infty}
   [A_n(x,\mathord{\cdot})B_n(\mathord{\cdot},y)]_m.
\]
It follows in particular that
\[
  A(x,\mathord{\cdot})(\mathord{\cdot},y)=A(x,y),
\]

\originalpage{180}
i.e. every bilinear form reproduces itself by convolution with $(x,y)$; we
may therefore also represent the value of the bilinear form as
\[
  A(x,y)
  =y_1\frac{\partial A}{\partial y_1}
   +y_2\frac{\partial A}{\partial y_2}+\cdots
  =x_1\frac{\partial A}{\partial x_1}
   +x_2\frac{\partial A}{\partial x_2}+\cdots,
\]
\end{lemma}

\begin{lemma}[4]\label{lem:4}
If $A,B,C,\ldots$ are bounded bilinear forms and the process of convolution
is repeatedly carried out with them, then the result is independent of the
order of the individual convolutions; for example,
\[
  \bigl(A(x,\mathord{\cdot})B(\mathord{\cdot},\mathord{\cdot})\bigr)
  C(\mathord{\cdot},y)
  =A(x,\mathord{\cdot})
  \bigl(B(\mathord{\cdot},\mathord{\cdot})C(\mathord{\cdot},y)\bigr).
\]
\end{lemma}

We now develop the simplest concepts and facts concerning orthogonal
transformations of infinitely many variables.  Let
\[
  o_{pq}\qquad(p,q=1,2,\ldots)
\]
be an arbitrary doubly indexed family of constants satisfying the relations
\[
  \sum_{q=1,2,\ldots}o_{pq}^{\,2}=1,
  \qquad
  \sum_{r=1,2,\ldots}o_{pr}o_{qr}=0\quad(p\ne q),
\]
and
\[
  \sum_{p=1,2,\ldots}o_{pq}^{\,2}=1,
  \qquad
  \sum_{r=1,2,\ldots}o_{rp}o_{rq}=0\quad(p\ne q).
\]
\begin{samepage}
Then the formulas
\begin{equation}
\tag{45}\label{eq:45}
\begin{aligned}
  x_1&=o_{11}x'_1+o_{12}x'_2+\cdots,\\
  x_2&=o_{21}x'_1+o_{22}x'_2+\cdots,\\
     &\ \vdots
\end{aligned}
\end{equation}
and
\[
\begin{aligned}
  x'_1&=o_{11}x_1+o_{21}x_2+\cdots,\\
  x'_2&=o_{12}x_1+o_{22}x_2+\cdots,\\
      &\ \vdots
\end{aligned}
\]
\end{samepage}
each define an orthogonal transformation; the latter transformation is the
inverse of the former.  The expressions on the right are bounded linear forms
in the infinitely many variables $x'_1,x'_2,\ldots$ and
$x_1,x_2,\ldots$, respectively.

The bilinear form
\[
  O(x,x')=\sum_{p,q=1,2,\ldots}o_{pq}x_p x'_q
\]
shall be called the bilinear form belonging to the orthogonal transformation
\eqref{eq:45}; as one readily sees from \originalpage{181}Fact I stated above,
it is certainly a bounded form.  By \hyperref[lem:3]{Lemma 3} one has
\begin{equation}
\tag{46}\label{eq:46}
  O(x,\mathord{\cdot})O(y,\mathord{\cdot})=(x,y);
\end{equation}
in particular,
\[
\begin{aligned}
  \sum_{p=1,2,\ldots}
  \bigl(o_{p1}x'_1+o_{p2}x'_2+\cdots\bigr)^2&=(x',x'),\\
  \sum_{p=1,2,\ldots}
  \bigl(o_{1p}x_1+o_{2p}x_2+\cdots\bigr)^2&=(x,x).
\end{aligned}
\]

If we apply the orthogonal transformation \eqref{eq:45} to any bounded
linear form $L(x)$, there results the linear form
\[
  L'(x')=L(\mathord{\cdot})O(\mathord{\cdot},x');
\]
consequently, if we subject both sets of variables of any bounded bilinear
form $A(x,y)$ to that transformation, there results the bilinear form
\[
  A'(x',y')
  =A(\mathord{\cdot},\mathord{\cdot})
   O(\mathord{\cdot},x')O(\mathord{\cdot},y').
\]
Let $A(x,y),B(x,y)$ be any two bounded bilinear forms, put
\[
  C(x,y)=A(x,\mathord{\cdot})B(\mathord{\cdot},y),
\]
and calculate the orthogonally transformed forms
\[
\begin{aligned}
  A'(x',y')
    &=A(\mathord{\cdot},\mathord{\cdot})
      O(\mathord{\cdot},x')O(\mathord{\cdot},y'),\\
  B'(x',y')
    &=B(\mathord{\circ},*)O(\mathord{\circ},x')O(*,y'),\\
  C'(x',y')
    &=C(\mathord{\cdot},\mathord{\cdot})
      O(\mathord{\cdot},x')O(\mathord{\cdot},y').
\end{aligned}
\]
As the convolution of the transformed forms we then find
\[
\begin{aligned}
  A'(x,\mathord{\circ})B'(\mathord{\circ},y)
  ={}&A(\mathord{\cdot},\mathord{\cdot})
      O(\mathord{\cdot},x)O(\mathord{\cdot},\mathord{\circ})\\
     &\quad{}B(*,*)O(*,\mathord{\circ})O(*,y),
\end{aligned}
\]
and, using \hyperref[lem:3]{Lemmas 3}, \hyperref[lem:4]{4}, and formula
\eqref{eq:46},
\[
\begin{aligned}
  A'(x,\mathord{\circ})B'(\mathord{\circ},y)
  &=A(\mathord{\cdot},\mathord{\circ})B(\mathord{\circ},*)
    O(\mathord{\cdot},x)O(*,y)\\
  &=C(\mathord{\cdot},*)O(\mathord{\cdot},x)O(*,y),
\end{aligned}
\]
and hence
\begin{equation}
\tag{47}\label{eq:47}
  A'(x,\mathord{\cdot})B'(\mathord{\cdot},y)=C'(x,y);
\end{equation}
i.e. the convolution of two bilinear forms is a covariant under an
orthogonal transformation.

If the sum of the coefficients of $x_p y_p$ in a bilinear form $A(x,y)$
converges, we denote this sum generally, as for finite forms, by
$A(\mathord{\cdot},\mathord{\cdot})$.

\originalpage{182}
We mention here the following fact, which is likewise easy to prove.

\begin{lemma}[5]\label{lem:5}
If
\[
  Q(x)=\sum_{p,q=1,2,\ldots}c_{pq}x_p x_q
\]
is a quadratic form of such a kind that
\[
  Q(\mathord{\cdot},\mathord{\cdot})
  Q(\mathord{\cdot},\mathord{\cdot})
  =\sum_{p,q=1,2,\ldots}c_{pq}^{\,2}
\]
converges to a finite limit, then this limit is an invariant under an
orthogonal transformation of $Q(x)$; i.e. that limit agrees with
\[
  Q'(\mathord{\cdot},\mathord{\cdot})
  Q'(\mathord{\cdot},\mathord{\cdot})
  =\sum_{p,q=1,2,\ldots}{c'_{pq}}^{\,2},
\]
where
\[
  Q'(x')=\sum_{p,q=1,2,\ldots}c'_{pq}x'_p x'_q
\]
denotes the quadratic form obtained from $Q(x)$ by orthogonal
transformation.
\end{lemma}

We now return to the theory developed above for the quadratic form $K(x)$,
and assume that this quadratic form $K(x)$ is a bounded form.

If, as before,
\[
  K_n(x)=[K(x)]_n
\]
denotes the $n$th truncation of $K(x)$, then the greatest value assumed by
$K_n(x)$ in absolute value is equal to
$1/\lvert\lambda_1^{(n)}\rvert$, where $\lambda_1^{(n)}$ denotes the
absolutely smallest of the $n$ eigenvalues of $K_n(x)$.  Since $K(x)$ is to
be a bounded form, there is a positive constant $M$ such that, for all $n$
and every system of values of the variables,
\[
  \lvert K_n(x)\rvert\le M,
\]
and hence also
\[
  \left\lvert\frac{1}{\lambda_1^{(n)}}\right\rvert\le M,
  \qquad\text{or}\qquad
  \lvert\lambda_1^{(n)}\rvert\ge\frac1M;
\]
i.e. the absolute values of the eigenvalues of $K_n(x)$ all remain above a
positive quantity different from zero, and thus $\lambda=0$ certainly does
not belong to the spectrum of $K(x)$.  Conversely, take a quadratic form
$K(x)$\originalpage{183} and suppose that $\lambda=0$ does not belong to its
spectrum.  Then, from some $n$ onward, all eigenvalues
$\lambda_p^{(n)}$ of its truncations $K_n(x)$ must remain, in absolute
value, above a positive bound $m$ different from zero.  It follows in turn
that the maxima $1/\lvert\lambda_1^{(n)}\rvert$ of the absolute values of the
truncations $K_n(x)$ must remain below the bound $1/m$; i.e. the form $K(x)$
is necessarily bounded.

Thus the assumption just made, and retained throughout what follows, that
$K(x)$ is a bounded form is completely equivalent to the statement that
$\lambda=0$ does not belong to the spectrum of $K(x)$, whereas the absolute
values of the eigenvalues of $K(x)$ may very well grow beyond every bound.

We now choose a quantity $\alpha$ so small that $\lambda=\alpha$ likewise
does not belong to the spectrum of $K(x)$.  Since the zeros of the
discriminants of the truncations of
\[
  (x,x)-\lambda K(x)
\]
then do not accumulate at $\lambda=\alpha$, those of the truncations of
\[
  (x,x)-\lambda\bigl\{(x,x)-\alpha K(x)\bigr\}
\]
do not, in absolute value, grow beyond every bound; i.e. the quadratic form
\[
  K^*(x)=(x,x)-\alpha K(x)
\]
does not have $\lambda=\infty$ as a condensation value; it is also bounded.
If $K_n(\lambda;x,y)$ denotes the resolvent of $K_n(x)$ and
$K_n^*(\lambda;x,y)$ that of $K_n^*(x)$, the definition of the resolvent
immediately gives
\begin{equation}
\tag{48}\label{eq:48}
  K_n(\lambda;x,y)
  =\frac{1}{1-\lambda/\alpha}
   K_n^*\!\left(\frac{\lambda}{\lambda-\alpha};x,y\right).
\end{equation}

We now apply \hyperref[thm:I]{Theorem I} to the form $K^*(x)$ and denote the
eigenvalues, eigenforms, interval spectrum, and spectral form of $K^*(x)$,
respectively, by
\[
  \lambda_1^*,\lambda_2^*,\ldots,
  \qquad E_1^*,E_2^*,\ldots,
  \qquad s^*,\sigma^*.
\]
Then, for all values of $\lambda^*$ lying outside the spectrum of $K^*$,
there results the equation, valid for every truncation,
\begin{equation}
\tag{49}\label{eq:49}
  \mathop{\mathrm{L}}\limits_{h=\infty}
  K_{m_h}^*(\lambda^*;x)
  =\sum_{(p)}\frac{E_p^*}{1-\lambda^*/\lambda_p^*}
   +\int_{(s^*)}\frac{d\sigma^*(\mu^*)}{1-\lambda^*/\mu^*},
\end{equation}
\originalpage{184}and likewise
\[
  (x,x)=\sum_{(p)}E_p^*+\int_{(s^*)}d\sigma^*(\lambda^*).
\]
In \eqref{eq:49} put
\[
  \lambda^*=\frac{\lambda}{\lambda-\alpha},
  \qquad
  \lambda_p^*=\frac{\lambda_p}{\lambda_p-\alpha},
  \qquad
  \mu^*=\frac{\mu}{\mu-\alpha},
\]
where
\[
  \frac{1}{1-\lambda^*/\lambda_p^*}
  =\frac{1-\lambda/\alpha}{1-\lambda/\lambda_p},
  \qquad
  \frac{1}{1-\lambda^*/\mu^*}
  =\frac{1-\lambda/\alpha}{1-\lambda/\mu}.
\]
If $s$ denotes the set of points $\mu$ corresponding to the set $s^*$ of
points $\mu^*$, then, in view of \eqref{eq:48}, we obtain the equation valid
for every truncation
\[
  \mathop{\mathrm{L}}\limits_{h=\infty}K_{m_h}(\lambda;x)
  =\sum_{(p,\infty)}\frac{E_p}{1-\lambda/\lambda_p}
   +\int_{(s)}\frac{d\sigma(\mu)}{1-\lambda/\mu}.
\]
Here the $\lambda_p$ are to be called the eigenvalues,
$E_p=E_p^*$ the corresponding eigenforms, $s$ the interval spectrum, and
$\sigma(\mu)=\sigma^*(\mu^*)$ the spectral form of $K(x)$; moreover,
$\lambda=\infty$ is to be included as an eigenvalue and $E_\infty$ as its
corresponding eigenfunction if $\lambda^*=1$ was an eigenvalue of $K^*(x)$.
From the formula above for $(x,x)$ we obtain
\begin{equation}
\tag{50}\label{eq:50}
  (x,x)=\sum_{(p,\infty)}E_p+\int_{(s)}d\sigma(\lambda).
\end{equation}

The last two formulas agree with \eqref{eq:42} and \eqref{eq:43} when
$K(x)$ does not have $\lambda=\infty$ as a condensation value.  Indeed, it
then follows from the definition of $K^*(x)$ that $\lambda^*=1$ is not a
condensation value, and therefore also not an eigenvalue, of $K^*$; hence
$\lambda=\infty$ is certainly not an eigenvalue of $K(x)$.

If $K(x,y)$ remains in absolute value below the finite bound $M$ for all
variables $x$, then from \eqref{eq:2} and \hyperref[lem:2]{Lemma 2} we infer
that, for every $m$,
\[
\begin{aligned}
  K_m(\lambda;x,y)
  ={}&(x,y)_m+\lambda K_m(x,y)+\lambda^2K_mK_m(x,y)+\cdots\\
     &+\lambda^nK_mK_m\cdots K_m(x,y)
      +\frac{\vartheta_n(\lambda M)^{n+1}}{1-\lambda M},
\end{aligned}
\]
\originalpage{185}where
\[
  n<m,\qquad \lvert\lambda\rvert<\frac1M,
  \qquad -1<\vartheta_n<+1.
\]
In view of \hyperref[lem:3]{Lemma 3}, if for $x,y$ we take fixed values all
of which vanish from some finite index onward, then, on putting $m=m_h$ and
passing to the limit $h\to\infty$, we obtain
\begin{equation}
\tag{51}\label{eq:51}
\begin{aligned}
  \mathop{\mathrm{L}}\limits_{h=\infty}K_{m_h}(\lambda;x,y)
  ={}&(x,y)+\lambda K(x,y)+\lambda^2KK(x,y)+\cdots\\
     &+\lambda^nKK\cdots K(x,y)
       +\frac{\vartheta_n(\lambda M)^{n+1}}{1-\lambda M},
\end{aligned}
\end{equation}
and from this, on letting $n\to\infty$,
\[
  \mathop{\mathrm{L}}\limits_{h=\infty}K_{m_h}(\lambda;x,y)
  =(x,y)+\lambda K(x,y)+\lambda^2KK(x,y)+\cdots.
\]
The formula thus obtained, together with the fact that every truncation of
the resolvent is regular analytic in $\lambda$ outside the spectrum, shows
that the resolvent
\[
  K(\lambda;x,y)
  =\mathop{\mathrm{L}}\limits_{h=\infty}K_{m_h}(\lambda;x,y)
\]
of the form $K(x,y)$ is determined uniquely by $K(x)$.  Returning to the
proof of the existence of the limit
\[
  \mathop{\mathrm{L}}\limits_{h=\infty}K_{m_h}(\lambda;x,y),
\]
it follows in turn that, in general, the limit
\[
  \mathop{\mathrm{L}}\limits_{m=\infty}K_m(\lambda;x,y)
\]
also exists and must equal the preceding limit.  The last formulas always
hold for every truncation of the forms under consideration.

From \eqref{eq:51} we conclude that the difference
\begin{equation}
\tag{52}\label{eq:52}
\begin{aligned}
  K(\lambda;x,y)-\bigl\{&(x,y)+\lambda K(x,y)+\lambda^2KK(x,y)+\cdots\\
  &{}+\lambda^nKK\cdots K(x,y)\bigr\}
\end{aligned}
\end{equation}
is a form for which, when $\lvert\lambda\rvert<1/M$, the absolute value of
every truncation remains below
\[
  \frac{\lvert\lambda M\rvert^{n+1}}{1-\lvert\lambda M\rvert}.
\]
By convolving the expression \eqref{eq:52} with $K(x,y)$, it follows from
\hyperref[lem:2]{Lemma 2} that the expression
\begin{equation}
\tag{53}\label{eq:53}
\begin{aligned}
  K(x,\mathord{\cdot})K(\lambda;\mathord{\cdot},y)
  -\bigl\{&K(x,y)+\lambda KK(x,y)+\lambda^2KKK(x,y)+\cdots\\
           &{}+\lambda^nKKK\cdots K(x,y)\bigr\}
\end{aligned}
\end{equation}

\originalpage{186}
remains in absolute value, in the same sense, below
\[
  M\frac{\lvert\lambda M\rvert^{n+1}}
          {1-\lvert\lambda M\rvert}.
\]
If in \eqref{eq:52} we replace $n$ by $n+1$ and then subtract from it
$\lambda$ times expression \eqref{eq:53}, there results the expression
\[
  K(\lambda;x,y)
  -\lambda K(x,\mathord{\cdot})K(\lambda;\mathord{\cdot},y)-(x,y),
\]
and this expression therefore remains in absolute value below
\[
  \frac{2\lvert\lambda M\rvert^{n+2}}
       {1-\lvert\lambda M\rvert}.
\]
Since this quantity tends to zero as $n\to\infty$, we have thereby proved
the validity of the equation
\begin{equation}
\tag{54}\label{eq:54}
  K(\lambda;x,y)
  -\lambda K(x,\mathord{\cdot})K(\lambda;\mathord{\cdot},y)
  =(x,y)
\end{equation}
for $\lvert\lambda\rvert<1/M$ and for every truncation, and consequently,
by a remark made above (\hyperref[origpage:177]{p.~177}), also for arbitrary
values of the infinitely many variables.

We arrived above at the equation
\begin{equation}
\tag{55}\label{eq:55}
  K(\lambda;x,y)
  =(x,y)+\lambda K(x,y)+\lambda^2KK(x,y)+\cdots
\end{equation}
and recognized its validity for every truncation and for
$\lvert\lambda\rvert<1/M$.  From this circumstance in turn we conclude that,
for arbitrary values of the variables and $\lvert\lambda\rvert<1/M$, every
truncation of
\begin{equation}
\tag{56}\label{eq:56}
\begin{aligned}
  K(\lambda;x,y)-\bigl\{&(x,y)+\lambda K(x,y)+\cdots\\
  &{}+\lambda^nKK\cdots K(x,y)\bigr\}
\end{aligned}
\end{equation}
is smaller in absolute value than
\begin{equation}
\tag{57}\label{eq:57}
  \frac{\lvert\lambda M\rvert^{n+1}}
       {1-\lvert\lambda M\rvert}.
\end{equation}
Hence, by a remark made above (\hyperref[origpage:177]{p.~177}), expression
\eqref{eq:56} itself must, for arbitrary values of the infinitely many
variables, remain in absolute value less than or equal to the quantity
\eqref{eq:57}.  Letting $n$ increase without limit, we conclude that equation
\eqref{eq:55} is valid for arbitrary values of the infinitely many variables
$x_1,x_2,\ldots$ and $y_1,y_2,\ldots$, and for
$\lvert\lambda\rvert<1/M$.

\originalpage{187}
The equation
\[
  K(\lambda,x)
  =\sum_{(p,\infty)}\frac{E_p}{1-\lambda/\lambda_p}
   +\int_{(s)}\frac{d\sigma(\mu)}{1-\lambda/\mu}
\]
was recognized above---as was \eqref{eq:50}---only as being valid in the
sense that the same truncations are understood to be taken on both sides.
We shall now show that this equation is valid for arbitrary values of the
infinitely many variables, provided that $\lambda$ denotes a value lying
outside the spectrum of $K$.

Let $a_1,a_2,\ldots$ be a fixed system of values of the infinitely many
variables $x_1,x_2,\ldots$; then let
\[
  E_p(a),\qquad [E_p(a)]_n,\qquad [E_p(a)]_m
\]
denote the values assumed, for that fixed system, by $E_p$ and by the $n$th
and $m$th truncations of $E_p$, respectively.  By \eqref{eq:50},
\[
  \sum_{(p=1,\ldots,P)}[E_p(a)]_n
\]
lies below a finite bound independent of $n$ and $P$.  We put
\[
  \sum_{(p=P+1,P+2,\ldots)}[E_p(a)]_n=\varepsilon(n,P),
\]
where $\varepsilon(n,P)$ is a quantity which, for fixed $n$, tends to zero as
$P\to\infty$.  Since, for fixed $\lambda$, the quantities
\[
  \left\lvert\frac{1}{1-\lambda/\lambda_p}\right\rvert
\]
all possess a finite upper bound $G$, the quantity $\eta(n,P)$ defined by
\[
  \sum_{(p=P+1,P+2,\ldots)}
  \frac{[E_p(a)]_n}{1-\lambda/\lambda_p}=\eta(n,P)
\]
likewise tends to zero as $P\to\infty$ for fixed $n$.

On the other hand,
\[
  \sum_{(p=1,\ldots,P)}\frac{E_p}{1-\lambda/\lambda_p}
\]
is a bounded form, and indeed one for which the absolute values of all its
values remain below a bound $G'$ independent of $P$.  It follows from our
earlier considerations\originalpage{188}
(\hyperref[origpage:178]{p.~178}) that
\[
\begin{aligned}
  &\sum_{(p=1,\ldots,P)}
    \frac{[E_p(a)]_m}{1-\lambda/\lambda_p}
  -\sum_{(p=1,\ldots,P)}
    \frac{[E_p(a)]_n}{1-\lambda/\lambda_p}\\
  &\hspace{5em}
    =(G)\sqrt{a_{n+1}^2+a_{n+2}^2+\cdots+a_m^2},
    \qquad(m>n),
\end{aligned}
\]
where $(G)$ denotes a quantity lying between finite bounds independent of
$n,m,P$.

From the last two equations we obtain
\[
\begin{aligned}
  \sum_{(p=1,\ldots,P)}
    \frac{[E_p(a)]_m}{1-\lambda/\lambda_p}
  ={}&\sum_{(p=1,2,\ldots)}
    \frac{[E_p(a)]_n}{1-\lambda/\lambda_p}\\
  &+(G)\sqrt{a_{n+1}^2+a_{n+2}^2+\cdots+a_m^2}
    +\eta(n,P).
\end{aligned}
\]
If here we first let $m\to\infty$, then $P\to\infty$, and finally
$n\to\infty$, we obtain
\[
  \sum_{(p=1,2,\ldots)}\frac{E_p(a)}{1-\lambda/\lambda_p}
  =\mathop{\mathrm{L}}\limits_{n=\infty}
   \sum_{(p=1,2,\ldots)}
   \frac{[E_p(a)]_n}{1-\lambda/\lambda_p}.
\]

On the other hand, we apply to the bilinear form $\sigma(\mu;x,y)$ the same
argument given above (\hyperref[origpage:178]{p.~178}) to prove the
convergence of $A(a,b)$.  Since, in view of \eqref{eq:31},
\[
  [\sigma(\mu;x,y)]_n\le1,
\]
and hence, for every $\mu$ and arbitrary $x_1,x_2,\ldots$ and
$y_1,y_2,\ldots$,
\[
  \sigma(\mu;x,y)\le1,
\]
that argument also proves that $[\sigma(\mu;x,y)]_n$ converges to
$\sigma(\mu;x,y)$ uniformly in $\mu$.  Thus, for every
system of values of the infinitely many variables $x_1,x_2,\ldots$ and
$y_1,y_2,\ldots$, the bilinear form $\sigma(\mu;x,y)$ represents a continuous
function of $\mu$, and it follows in turn that
\[
  \int_{(s)}\frac{d\sigma(\mu)}{1-\lambda/\mu}
  =\mathop{\mathrm{L}}\limits_{n=\infty}
   \left[\int_{(s)}\frac{d\sigma(\mu)}{1-\lambda/\mu}\right]_n.
\]
The last two limit equations prove our assertion, i.e. extend the validity of
the partial-fraction representation of $K(\lambda;x)$ to arbitrary values of
the infinitely many variables.

The partial-fraction representation of $K(\lambda;x)$ just proved to hold in
general shows at the same time that the value of the resolvent
\originalpage{189}$K(\lambda;x)$, for any system of the infinitely
many variables $x_1,x_2,\ldots$, is a regular analytic function of $\lambda$
outside the spectrum.

If, in $K(\lambda;x,y)$, we substitute generally for $x_p$ the expression
$\partial K(x,y)/\partial y_p$, there results the convolution
\[
  K(x,\mathord{\cdot})K(\lambda;\mathord{\cdot},y),
\]
and hence this too, for every system of values of the infinitely many
variables $x_1,x_2,\ldots$ and $y_1,y_2,\ldots$, represents a regular analytic
function of $\lambda$ outside the spectrum.

From these facts we infer the validity of equation \eqref{eq:54}, not only
for arbitrary systems of values of the infinitely many variables
$x_1,x_2,\ldots$ and $y_1,y_2,\ldots$, but also for all values of $\lambda$
lying outside the spectrum.

We summarize the most important of the results obtained as follows.

\begin{theorem}[II]\label{thm:II}
Let $K(x)$ be a bounded quadratic form in the infinitely many variables
$x_1,x_2,\ldots$.  The resolvent $K(\lambda,x)$ of $K(x)$ is a uniquely
determined quadratic form in these same variables,
\[
  K(\lambda,x)=\sum_{(p,q)}K_{pq}(\lambda)x_p x_q,
\]
whose coefficients $K_{pq}(\lambda)$ are regular analytic functions of
$\lambda$ for all values of $\lambda$ lying outside the spectrum of $K$.

When $\lambda$ denotes a value lying outside the spectrum of $K$, the
resolvent $K(\lambda,x)$ is a bounded form; for any system of
values of the infinitely many variables $x_1,x_2,\ldots$, it represents an
analytic function of $\lambda$.

For arbitrary values of the infinitely many variables $x_1,x_2,\ldots$ and
for sufficiently small values of $\lambda$, the resolvent $K(\lambda,x)$
admits the power-series expansion
\begin{equation}
\tag{59}\label{eq:59}
  K(\lambda,x)=(x,x)+\lambda K(x)+\lambda^2KK(x)+\cdots,
\end{equation}
and, likewise for arbitrary values of the infinitely many variables and for
all values of $\lambda$ lying outside the spectrum of $K$, the
partial-fraction representation
\begin{equation}
\tag{60}\label{eq:60}
  K(\lambda,x)
  =\sum_{(p,\infty)}\frac{E_p}{1-\lambda/\lambda_p}
   +\int_{(s)}\frac{d\sigma(\mu)}{1-\lambda/\mu}.
\end{equation}
Here the sum is to be \originalpage{190}extended over the entire point
spectrum of $K$, i.e. over all eigenvalues, possibly including the eigenvalue
$\infty$; $E_p$ denotes in general the quadratic eigenform belonging
to $\lambda_p$.  It is a bounded form which is nonnegative for every system
of values of the variables $x_1,x_2,\ldots$.  The integral is to be extended
over the interval spectrum of $K$.  The spectral form $\sigma(\lambda)$ is a
bounded form in the infinitely many variables $x_1,x_2,\ldots$; for each
system of their values it represents, as a function of $\lambda$, a function
which is continuous and, within the interval spectrum $s$, increases with
$\lambda$---apart from special values of $x_1,x_2,\ldots$---but remains
constant in every interval lying outside $s$.

In particular, the equations
\begin{equation}
\tag{61}\label{eq:61}
  (x,x)=\sum_{(p,\infty)}E_p+\int_{(s)}d\sigma(\lambda),
\end{equation}
\begin{equation}
\tag{62}\label{eq:62}
  K(x)=\sum_{(p)}\frac{E_p}{\lambda_p}
       +\int_{(s)}\frac{d\sigma(\mu)}{\mu}
\end{equation}
hold.  The resolvent $K(\lambda,x)$ is connected with $K(x)$ by the relation
\begin{equation}
\tag{63}\label{eq:63}
  K(\lambda;x,y)
  -\lambda K(x,\mathord{\cdot})K(\lambda;\mathord{\cdot},y)
  =(x,y),
\end{equation}
which is valid for all values of $\lambda$ lying outside the spectrum of
$K$.
\end{theorem}

Put
\[
  K(\lambda;x,y)=\alpha_1x_1+\alpha_2x_2+\cdots,
\]
where $\alpha_1,\alpha_2,\ldots$ are certain linear functions of
$y_1,y_2,\ldots$ with convergent sum of squares.  Then, by equating the
coefficients of $x_p$ in \eqref{eq:63}, we obtain
\[
  \alpha_p-\lambda\sum_{(q)}k_{pq}\alpha_q=y_p,
  \qquad(p=1,2,\ldots).
\]
Thus $\alpha_1,\alpha_2,\ldots$ solve these infinitely many inhomogeneous
equations arising from the quadratic form $(x,x)-\lambda K(x)$ (where
$\lambda$ lies outside the spectrum), whenever $y_1,y_2,\ldots$ are any
quantities with convergent sum of squares; they are the unique solution with
convergent sum of squares.

We shall now investigate the behavior of the resolvent $K(\lambda,x)$ for a
value of $\lambda$ lying within the spectrum of $K$.

For this purpose put
\[
  \lambda=\nu+i\nu',
\]
where $\nu,\nu'$ are real numbers, and ask whether the product
\[
  (\nu-\lambda)
  \left\{
    \sum_{(p)}\frac{E_p}{\lambda_p-\lambda}
    +\int_{(s)}\frac{d\sigma(\mu)}{\mu-\lambda}
  \right\}
\]
tends to a limit as $\nu'\to0$.

\originalpage{191}
First let $\nu$ be a condensation point of the eigenvalues of $K(x)$, but not
equal to an eigenvalue of $K(x)$.  Put
\[
  \sum_{(p)}E_p=E_1+\cdots+E_m+R_m,
\]
and denote by $\nu_m$ that one among the eigenvalues
$\lambda_1,\ldots,\lambda_m$ which lies nearest to $\nu$.  Choose $m$ so
large that
\[
  \nu'\le(\nu-\nu_m)^2
\]
still holds.
Then
\[
  \text{for }p\le m,\qquad
  \left\lvert\frac{\nu-\lambda}{\lambda_p-\lambda}\right\rvert
  \le\left\lvert\frac{\nu'}{\lambda_p-\nu}\right\rvert
  \le\frac{\nu'}{\lvert\nu_m-\nu\rvert}
  \le\sqrt{\nu'},
\]
and, for $p\ge m$, certainly
\[
  \left\lvert\frac{\nu-\lambda}{\lambda_p-\lambda}\right\rvert\le1;
\]
hence also
\[
  \left\lvert
    (\nu-\lambda)\sum_{(p)}\frac{E_p}{\lambda_p-\lambda}
  \right\rvert
  \le\sqrt{\nu'}(E_1+\cdots+E_m)+R_m
  \le\sqrt{\nu'}+R_m.
\]
Since, as $\nu'\to0$, $m$ necessarily grows beyond every bound and hence
$R_m$ tends to zero, it follows that
\begin{equation}
\tag{64}\label{eq:64}
  \mathop{\mathrm{L}}\limits_{\nu'=0}
  \left\{(\nu-\lambda)
  \sum_{(p)}\frac{E_p}{\lambda_p-\lambda}\right\}=0.
\end{equation}
The last limit equation certainly also holds if $\nu$ is neither a
condensation point of the eigenvalues of $K(x)$ nor itself equal to an
eigenvalue.

From these facts we infer, on the other hand, that if $\nu$ equals the
eigenvalue $\lambda_p$, then necessarily
\begin{equation}
\tag{65}\label{eq:65}
  \mathop{\mathrm{L}}\limits_{\nu'=0}
  \left\{(\lambda_p-\lambda)
  \sum_{(p)}\frac{E_p}{\lambda_p-\lambda}\right\}=E_p.
\end{equation}

Now let $\nu$ be a point of the interval spectrum $s$, and let $\mu'$ be a
real number greater than $\nu$.  If we put
\[
  \nu'=(\nu-\mu')^2,
\]
then, by a similar estimate of the integral extending up to $\mu'$, we obtain
\begin{equation}
\tag{66}\label{eq:66}
  \mathop{\mathrm{L}}\limits_{\nu'=0}
  \left\{(\nu-\lambda)
  \int_{\sigma(\nu)}^{\sigma(+\infty)}
  \frac{d\sigma(\mu)}{\mu-\lambda}\right\}=0.
\end{equation}
Likewise,
\begin{equation}
\tag{67}\label{eq:67}
  \mathop{\mathrm{L}}\limits_{\nu'=0}
  \left\{(\nu-\lambda)
  \int_{\sigma(-\infty)}^{\sigma(\nu)}
  \frac{d\sigma(\mu)}{\mu-\lambda}\right\}=0.
\end{equation}

\originalpage{192}
Since
\[
  K(\lambda,x)=(x,x)+\lambda
  \left\{
    \sum_{(p,\infty)}\frac{E_p}{\lambda_p-\lambda}
    +\int_{(s)}\frac{d\sigma(\mu)}{\mu-\lambda}
  \right\},
\]
it follows from \eqref{eq:64}, \eqref{eq:65}, \eqref{eq:66}, and
\eqref{eq:67} that
\begin{equation}
\tag{68}\label{eq:68}
  \mathop{\mathrm{L}}\limits_{\nu'=0}
  (\nu-\lambda)K(\lambda,x)
  =\begin{cases}
     0, & \text{if $\nu$ is none of the eigenvalues of $K(x)$},\\
     \lambda_pE_p, & \text{if $\nu$ equals the eigenvalue $\lambda_p$}.
   \end{cases}
\end{equation}
As a supplement to this, if
\[
  \lambda=i\nu',
\]
we have the limit equation
\begin{equation}
\tag{69}\label{eq:69}
  \mathop{\mathrm{L}}\limits_{\nu'=\infty}K(\lambda,x)
  =\begin{cases}
     0, & \text{if $\lambda=\infty$ is not an eigenvalue},\\
     E_\infty, & \text{if it must be counted as one}.
   \end{cases}
\end{equation}
To see this, the analogous considerations to those above suffice: in
examining the sum, choose $m$ so large that $\nu'>\nu_m^2$ still holds, where
$\nu_m$ denotes the eigenvalue of greatest absolute value among
$\lambda_1,\ldots,\lambda_m$, and in examining the integral put
$\nu'=\mu'^2$.

If $\mu$ denotes any value lying outside the spectrum of $K$, then
convolution of \eqref{eq:63} with $K(\mu;x,y)$ gives
\[
\begin{aligned}
  &\lambda K(\mathord{\cdot},\mathord{\cdot})
  K(\lambda;\mathord{\cdot},x)K(\mu;\mathord{\cdot},y)\\
  &\qquad
  =K(\lambda;\mathord{\cdot},x)K(\mu;\mathord{\cdot},y)-K(\mu;x,y).
\end{aligned}
\]
From this formula and the one obtained from it by simultaneously
interchanging $\lambda;x_1,x_2,\ldots$ with
$\mu;y_1,y_2,\ldots$, we find
\begin{equation}
\tag{70}\label{eq:70}
  (\lambda-\mu)K(\lambda;x,\mathord{\cdot})
  K(\mu;y,\mathord{\cdot})
  =\lambda K(\lambda;x,y)-\mu K(\mu;x,y).
\end{equation}

Now recall that the convolution
\[
  K(\lambda;x,\mathord{\cdot})K(\mu;y,\mathord{\cdot})
\]
means the expression obtained when, in $K(\lambda;x,y)$, we substitute for
the variables $y_1,y_2,\ldots$, respectively, the values
\[
  \frac{\partial K(\mu;x,y)}{\partial y_1},
  \quad
  \frac{\partial K(\mu;x,y)}{\partial y_2},\ldots.
\]
It follows from \eqref{eq:68} that, with $\mu$ held fixed and
\originalpage{193}for $\lambda=\lambda_p+i\nu'$, the limit equation
\begin{equation}
\tag{71}\label{eq:71}
  \mathop{\mathrm{L}}\limits_{\nu'=0}
  (\lambda_p-\lambda)
  K(\lambda;x,\mathord{\cdot})K(\mu;y,\mathord{\cdot})
  =\lambda_pE_p(x,\mathord{\cdot})K(\mu;y,\mathord{\cdot})
\end{equation}
must hold.  On the other hand, again in view of \eqref{eq:68}, in the same
sense
\begin{equation}
\tag{72}\label{eq:72}
\begin{aligned}
  \mathop{\mathrm{L}}\limits_{\nu'=0}
  (\lambda_p-\lambda)
  \bigl\{\lambda K(\lambda;x,y)-\mu K(\mu;x,y)\bigr\}
  =\lambda_p^2E_p(x,y).
\end{aligned}
\end{equation}
It therefore follows from \eqref{eq:71}, \eqref{eq:72}, and
\eqref{eq:70}, writing $\lambda$ in place of $\mu$, that
\begin{equation}
\tag{73}\label{eq:73}
  E_p(x,\mathord{\cdot})K(\lambda;y,\mathord{\cdot})
  =\frac{\lambda_p}{\lambda_p-\lambda}E_p(x,y).
\end{equation}
If in this equation we again put $\lambda=\lambda_p+i\nu'$, then
multiplication by $\lambda_p-\lambda$ and application of \eqref{eq:68} give
\begin{equation}
\tag{74}\label{eq:74}
  E_p(x,\mathord{\cdot})E_p(\mathord{\cdot},y)=E_p(x,y).
\end{equation}
If, however, in \eqref{eq:73} we first replace $p$ by an index $q$ different
from $p$ and then proceed in the same way, we obtain
\begin{equation}
\tag{75}\label{eq:75}
  E_p(x,\mathord{\cdot})E_q(\mathord{\cdot},y)=0,
  \qquad(p\ne q).
\end{equation}
Using \eqref{eq:69}, we see in the same way that \eqref{eq:74} and
\eqref{eq:75} are also valid when the eigenform $E_\infty$ that may belong to
$\lambda=\infty$ is put in place of $E_p$.

If the value of $K(\lambda;x,y)$ from \eqref{eq:60} is substituted in
\eqref{eq:73}, then \eqref{eq:74} and \eqref{eq:75} show that, identically in
$\lambda$,
\begin{equation}
\tag{76}\label{eq:76}
  E_p(x,\mathord{\cdot})
  \int_{(s)}\frac{d\sigma(\mu;\mathord{\cdot},y)}{1-\lambda/\mu}=0;
\end{equation}
the same equation also holds, when applicable, for the eigenform $E_\infty$.

In the discussion that follows, by an \emph{elementary form} we understand in
general a bounded quadratic form $E$ whose point spectrum in the finite plane
consists only of the single point $1$ and which has no interval spectrum.
Applying representation \eqref{eq:62} to the elementary form $E$, we find
that $E$ itself is the eigenform belonging to the eigenvalue $1$, and hence,
    by \eqref{eq:74}, \originalpage{194}it must satisfy
\begin{equation}
\tag{77}\label{eq:77}
  E(x,\mathord{\cdot})E(\mathord{\cdot},y)=E(x,y).
\end{equation}
Conversely, if a bounded quadratic form $E$ satisfies relation
\eqref{eq:77}, then application of formula \eqref{eq:59} gives for its
resolvent the expression
\[
  (x,x)+\lambda E+\lambda^2E+\cdots
  =(x,x)-E+\frac{E}{1-\lambda}.
\]
From this, in view of \eqref{eq:68} and \eqref{eq:69}, we see that $E$ has
only the one finite eigenvalue $1$; it follows further that there is no
interval spectrum.  Thus $E$ is an elementary form.

Let
\[
\begin{aligned}
  L_1(x)&=l_{11}x_1+l_{12}x_2+\cdots,\\
  L_2(x)&=l_{21}x_1+l_{22}x_2+\cdots,\\
        &\ \vdots
\end{aligned}
\]
be any finite or infinite number of linear forms whose coefficients satisfy
\[
  L_p(\mathord{\cdot})L_p(\mathord{\cdot})
  =\sum_{(r)}l_{pr}^{\,2}=1,
  \qquad
  L_p(\mathord{\cdot})L_q(\mathord{\cdot})
  =\sum_{(r)}l_{pr}l_{qr}=0.
\]
Such a system of forms shall be called a \emph{system of orthogonal linear
forms}, or briefly an \emph{orthogonal system}.

The close connection between the concept just defined and that of an
elementary form is exhibited by the following theorem.

\begin{theorem}
Every elementary form can be represented as a sum of squares of the linear
forms of an orthogonal system; conversely, the sum of the squares of the
linear forms of an orthogonal system always represents an elementary form.
\end{theorem}

\begin{proof}
For the first assertion, observe that the elementary form $E$, since it is
itself its eigenform, is necessarily nonnegative.  Thus, if the variable
$x_1$ occurs at all in $E$, then the coefficient of $x_1^2$ in $E$---denote
it by $e_{11}$---is certainly positive.  Put
\[
  L_1(x)=\frac{1}{\sqrt{e_{11}}}
  \frac{\partial E(x,y)}{\partial y_1}.
\]
Then \eqref{eq:77} gives
\[
  L_1(\mathord{\cdot})E(x,\mathord{\cdot})=L_1(x),
  \qquad
  L_1(\mathord{\cdot})L_1(\mathord{\cdot})=1.
\]
Consequently, if we form
\begin{equation}
\tag{78}\label{eq:78}
  E_1(x)=E(x)-L_1^2(x),
\end{equation}
\originalpage{195}we obtain
\begin{equation}
\tag{79}\label{eq:79}
\begin{aligned}
  L_1(\mathord{\cdot})E_1(x,\mathord{\cdot})&=0,\\
  E_1(x,\mathord{\cdot})E_1(\mathord{\cdot},y)&=E_1(x,y).
\end{aligned}
\end{equation}
Thus $E_1$ is likewise an elementary form.  Since, as such, $E_1$ is a
nonnegative form, and since by \eqref{eq:78} the coefficient of $x_1^2$ in
$E_1$ has the value zero, the variable $x_1$ does not occur at all in $E_1$.

Applying the same procedure to $E_1$ instead of $E$, we arrive at a linear
form $L_2$ and at the elementary form
\[
  E_2(x)=E_1(x)-L_2^2(x)=E(x)-L_1^2(x)-L_2^2(x),
\]
which does not contain the variables $x_1,x_2$; at the same time,
\eqref{eq:79} gives
\[
  L_1(\mathord{\cdot})L_2(\mathord{\cdot})=0.
\]
Finally, the expression
\[
  E(x)-L_1^2(x)-L_2^2(x)-\cdots
\]
is a bounded nonnegative form which contains none of the variables
$x_1,x_2,\ldots$ and is therefore identically zero; i.e.
\[
  E(x)=L_1^2(x)+L_2^2(x)+\cdots.
\]

To prove the converse assertion of the theorem, first form, from the first
$m$ linear forms $L_1,\ldots,L_m$ of the given orthogonal system, the
expression linear in $x_1,x_2,\ldots$
\[
  M(x)=(x,y)-L_1(x)L_1(y)-\cdots-L_m(x)L_m(y).
\]
Since the sum of the squares of the coefficients of $M(x)$ is
\[
  M(\mathord{\cdot})M(\mathord{\cdot})
  =(y,y)-L_1^2(y)-\cdots-L_m^2(y),
\]
the right-hand side here is nonnegative.  Consequently, the finite, or
infinitely continued, series of forms
\[
  L_1^2(y)+L_2^2(y)+\cdots
\]
also represents a bounded quadratic form; and from the orthogonality
properties of the linear forms $L_1,L_2,\ldots$ it then follows that this
form satisfies relation \eqref{eq:77}.
\end{proof}

The elementary form $(x,x)$, and only this one, also does not have
$\lambda=\infty$ as an eigenvalue.

We now apply the preceding results to the eigenforms of the quadratic form
$K(x)$.  If we regard all the eigenforms $E_p$, or $E_p,E_\infty$, of $K(x)$,
which are elementary forms by \eqref{eq:74}, as represented by sums of
squares of orthogonal linear\originalpage{196} forms
$x'_1,x'_2,\ldots$, we arrive in the following way at an orthogonal
substitution of the variables $x_1,x_2,\ldots$.

First form
\[
  (x,x)-{x'_1}^{\,2}-{x'_2}^{\,2}-\cdots;
\]
it satisfies relation \eqref{eq:77} and is therefore an elementary form in
the variables $x_1,x_2,\ldots$.  If we put it in the form
\[
  \xi_1^2+\xi_2^2+\cdots,
\]
where $\xi_1,\xi_2,\ldots$ likewise denote an orthogonal system of linear
forms which are also orthogonal to $x'_1,x'_2,\ldots$, then
\[
  (x,x)={x'_1}^{\,2}+{x'_2}^{\,2}+\cdots+\xi_1^2+\xi_2^2+\cdots,
\]
and hence the linear forms $x'_1,x'_2,\ldots,\xi_1,\xi_2,\ldots$, taken
together, define an orthogonal substitution of the variables
$x_1,x_2,\ldots$.

Since equation \eqref{eq:76} certainly holds for all values of $\lambda$ in
a sufficiently small neighborhood of $\lambda=0$, and since every continuous
function $w(\mu)$ on the interval $s$, which does not contain $0$, can be
uniformly approximated by linear aggregates of the functions\footnote{Translator's
note: The original edition prints $1-\lambda/\mu$ at this point; equation
\eqref{eq:76} and the ensuing approximation argument require the reciprocal
shown here.}
\[
  \frac{1}{1-\lambda/\mu}\mkern6mu,
\]
we conclude that, for every continuous function $w(\mu)$,
\[
  E_p(x,\mathord{\cdot})
  \int_{(s)}w(\mu)\,d\sigma(\mu;\mathord{\cdot},y)=0,
\]
and consequently, for every $\mu$,
\begin{equation}
\tag{80}\label{eq:80}
  E_p(x,\mathord{\cdot})\sigma(\mu;\mathord{\cdot},y)=0.
\end{equation}
If in this equation we imagine the variables
$x'_1,x'_2,\ldots,\xi_1,\xi_2,\ldots$ introduced in place of
$x_1,x_2,\ldots$, then \eqref{eq:80}, in view of the covariance of
convolution under orthogonal transformation and the nonnegative character of
$\sigma(\mu;x)$, shows that the spectral form so transformed is free of the
variables $x'_1,x'_2,\ldots$; we denote it by $\xi(\mu;\xi)$.  Formula
\eqref{eq:60} then takes the form
\begin{equation}
\tag{81}\label{eq:81}
  K(\lambda,x)
  =\sum_{(k=1,2,\ldots)}\frac{{x'_k}^{\,2}}{1-k_k\lambda}
   +\int_{(s)}\frac{d\xi(\mu;\xi)}{1-\lambda/\mu}\mkern6mu,
\end{equation}
where $k_1,k_2,\ldots$ denote the corresponding reciprocal eigenvalues,
after they have been arranged in a simply infinite sequence.

To find the characteristic properties of the spectral form $\xi(\mu;\xi)$,
we substitute in \eqref{eq:70} \originalpage{197}the expression
\eqref{eq:81} just found for the resolvent. We obtain
\[
\begin{aligned}
  &(\lambda-\mu)
  \int_{(s)}\frac{d\xi(\rho;\xi,\mathord{\cdot})}{1-\lambda/\rho}
  \int_{(s)}\frac{d\xi(\rho;\eta,\mathord{\cdot})}{1-\mu/\rho}\\
  &\qquad
  =\lambda\int_{(s)}\frac{d\xi(\rho;\xi,\eta)}{1-\lambda/\rho}
   -\mu\int_{(s)}\frac{d\xi(\rho;\xi,\eta)}{1-\mu/\rho}.
\end{aligned}
\]
In view of the identity
\[
  \frac{
    \dfrac{\lambda}{1-\lambda/\rho}
    -\dfrac{\mu}{1-\mu/\rho}
  }{\lambda-\mu}
  =\frac{1}{1-\lambda/\rho}\frac{1}{1-\mu/\rho},
\]
it follows that
\[
\begin{aligned}
  &\int_{(s)}\frac{d\xi(\rho;\xi,\mathord{\cdot})}{1-\lambda/\rho}
   \int_{(s)}\frac{d\xi(\rho;\eta,\mathord{\cdot})}{1-\mu/\rho}\\
  &\qquad
   =\int_{(s)}
    \frac{1}{1-\lambda/\rho}\frac{1}{1-\mu/\rho}
    \,d\xi(\rho;\xi,\eta).
\end{aligned}
\]
Since this equation is to hold for all values of $\lambda,\mu$, we conclude,
as above, that for arbitrary continuous functions $u(\rho),v(\rho)$,
\[
\begin{aligned}
  &\int_{(s)}u(\rho)\,d\xi(\rho;\xi,\mathord{\cdot})
   \int_{(s)}v(\rho)\,d\xi(\rho;\eta,\mathord{\cdot})\\
  &\qquad
   =\int_{(s)}u(\rho)v(\rho)\,d\xi(\rho;\xi,\eta).
\end{aligned}
\]
For $u(\rho)=v(\rho)$ and $\xi=\eta$, this relation takes the form
\[
  \int_{(s)}\bigl(u(\rho)\bigr)^2\,d\xi(\rho;\xi)
  =\int_{(s)}u(\rho)\,d\xi(\rho;\xi,\mathord{\cdot})
   \int_{(s)}u(\rho)\,d\xi(\rho;\xi,\mathord{\cdot}).
\]
Furthermore, \eqref{eq:81} with $\lambda=0$ gives
\[
  (\xi,\xi)=\int_{(s)}d\xi(\mu;\xi).
\]

Conversely, the preceding relations, and hence ultimately \eqref{eq:70}, can
be inferred from the penultimate relation, while \eqref{eq:63} can then be
inferred from \eqref{eq:70} with the aid of the last equation.  Thus we see
that the two conditions just found are also sufficient to characterize the
spectral form $\xi$.  From \eqref{eq:68}, \eqref{eq:69}, and \eqref{eq:74}
follows the unique determination of $\lambda_p,E_p$; and, on the other hand,
we conclude that the spectral form is also uniquely determined by the above
relations characteristic of it and by the requirement that $K$ admit
representation \eqref{eq:62}.  For the mode of reasoning just indicated,
starting from the properties of the spectral form,
\originalpage{198}shows that the resolvent is represented by \eqref{eq:60},
and hence, because of its unique determination, that
\[
  \int_{(s)}\frac{d\sigma(\mu)}{1-\lambda/\mu}
\]
is uniquely determined for every $\lambda$, and therefore so is the spectral
form itself.

We summarize the results obtained as follows.

\begin{theorem}[III]\label{thm:III}
Every bounded quadratic form $K$ in infinitely many variables can always,
and in only one way, be brought by an orthogonal substitution into the form
\[
  K=k_1x_1^2+k_2x_2^2+\cdots
    +\int_{(s)}\frac{d\sigma(\mu;\xi)}{\mu},
\]
where $k_1,k_2,\ldots$ denote certain quantities lying in absolute value
below a finite bound (the reciprocal eigenvalues), and the spectral form
$\sigma(\mu;\xi)$ satisfies the relations sufficient for its characterization
\begin{equation}
\tag{82}\label{eq:82}
\begin{aligned}
  \int_{(s)}\bigl(u(\mu)\bigr)^2\,d\sigma(\mu;\xi)
    &=\int_{(s)}u(\mu)\,d\sigma(\mu;\xi,\mathord{\cdot})
      \int_{(s)}u(\mu)\,d\sigma(\mu;\xi,\mathord{\cdot}),\\
  (\xi,\xi)&=\int_{(s)}d\sigma(\mu;\xi),
\end{aligned}
\end{equation}
identically for all continuous functions $u(\mu)$.
\end{theorem}

Put $\lambda=\lambda_p+i\nu'$ in \eqref{eq:63}.  After multiplication by
$\lambda_p-\lambda$ and taking the limit $\nu'\to0$, it follows from
\eqref{eq:68} that
\[
  E_p(x,y)-\lambda_pK(x,\mathord{\cdot})
  E_p(\mathord{\cdot},y)=0.
\]
If we represent $E_p$ as a sum of squares of an orthogonal system,
\[
  E_p=\sum_{(h)}\bigl(L_{ph}(x)\bigr)^2,
\]
then convolution with $L_{ph}(y)$ gives
\[
  L_{ph}(x)-\lambda_pK(x,\mathord{\cdot})
  L_{ph}(\mathord{\cdot})=0.
\]
Thus the equation
\[
  \lambda L(\mathord{\cdot})K(x,\mathord{\cdot})=L(x)
\]
can certainly be satisfied identically in $x_1,x_2,\ldots$ by a bounded
linear form $L$ whenever $\lambda$ is one of the eigenvalues of $K$.
Likewise, it follows from \eqref{eq:69} that
\[
  L(\mathord{\cdot})K(x,\mathord{\cdot})=0
\]
can certainly be satisfied when $\lambda=\infty$ is an eigenvalue.

We can now also establish conversely that the preceding equation is solvable
by a bounded linear form only in this case.  In view of
\hyperref[thm:III]{Theorem III}, it is necessary only
\originalpage{199}to prove that, if $L$ is a bounded linear form in the
variables $\xi_1,\xi_2,\ldots$, the relation
\[
  \lambda L(\mathord{\cdot})
  \int_{(s)}\frac{d\sigma(\mu;\xi,\mathord{\cdot})}{\mu}=L(\xi)
\]
cannot hold for any value of $\lambda$.  Denote the coefficients of $L$ by
$l$; the latter relation then takes the form
\[
  \lambda\int_{(s)}\frac{d\sigma(\mu;\xi,l)}{\mu}=(\xi,l).
\]
From this we subtract the relation
\[
  \int_{(s)}d\sigma(\mu;\xi,l)=(\xi,l)
\]
and obtain
\[
  \int_{(s)}\left(1-\frac{\lambda}{\mu}\right)
  d\sigma(\mu;\xi,l)=0,
\]
which would likewise be satisfied identically in $\xi_1,\xi_2,\ldots$.
Taking
\[
  u(\mu)=1-\frac{\lambda}{\mu}
\]
in \eqref{eq:82}, we infer that
\[
  \int_{(s)}\left(1-\frac{\lambda}{\mu}\right)^2
  d\sigma(\mu;l)=0.
\]
In view of
\[
  \int_{(s)}d\sigma(\mu;l)=(l,l),
\]
and because $\sigma$ never decreases, this is possible only when
$(l,l)=0$.  We thus obtain the following fact.

\begin{theorem}[IV]\label{thm:IV}
If $K$ is any bounded quadratic form, then the relation
\[
  \lambda L(\mathord{\cdot})K(x,\mathord{\cdot})=L(x)
\]
is solvable by a bounded linear form $L$ if and only if $\lambda$ is an
eigenvalue of $K$.
\end{theorem}

In particular, the equation
\[
  L(\mathord{\cdot})K(x,\mathord{\cdot})=0
\]
is solvable if and only if $\lambda=\infty$ is an eigenvalue of $K$.  If
$\lambda=\infty$ is not an eigenvalue, and this equation is therefore not
solvable, the quadratic form shall be called \emph{closed}.

\originalpage{200}
In the remainder of this \hyperref[sec:XI]{Section XI} we shall concern
ourselves in greater detail with two particular cases of
\hyperref[thm:III]{Theorem III}.

We call a function $F(x_1,x_2,\ldots)$ of infinitely many variables
$x_1,x_2,\ldots$ \emph{completely continuous} at a fixed system of their
values if the values of $F(x_1+\varepsilon_1,x_2+\varepsilon_2,\ldots)$
converge to $F(x_1,x_2,\ldots)$ whenever each $\varepsilon_j$ tends to zero,
independently of how this happens; that is, if
\[
  \mathop{\mathrm{L}}\limits_{\varepsilon_1=0,\,
    \varepsilon_2=0,\ldots}
  F(x_1+\varepsilon_1,x_2+\varepsilon_2,\ldots)
  =F(x_1,x_2,\ldots)
\]
for every sequence of systems
$\varepsilon_1^{(h)},\varepsilon_2^{(h)},\ldots$ whose individual coordinates
tend to zero, so that
\[
  \mathop{\mathrm{L}}\limits_{h=\infty}\varepsilon_1^{(h)}=0,
  \qquad
  \mathop{\mathrm{L}}\limits_{h=\infty}\varepsilon_2^{(h)}=0,
  \quad\ldots.
\]
If a function is completely continuous for every system of values of the
variables, it shall simply be called \emph{completely continuous}.  The
following conclusions attach immediately to the concept of complete
continuity.

If a completely continuous function $F$ has the property of remaining in
absolute value below a finite quantity for all values of the variables, then
$F$ possesses a maximum, as can readily be proved by applying the familiar
procedure for a finite number of variables.  If, furthermore,
$L_1(x),\ldots,L_m(x)$ denote $m$ further completely continuous functions of
the variables $x_1,x_2,\ldots$, and only those systems of values of these
variables are admitted which satisfy
\[
  L_1(x)=0,\quad\ldots,\quad L_m(x)=0,
\]
then $F$ possesses a relative maximum; here the variables are always subject
to the inequality
\[
  (x,x)\le1.
\]

A bounded linear form in the variables $x_1,x_2,\ldots$ is, as is seen at
once, also completely continuous in these variables.  From this we readily
conclude that a completely continuous function of the variables
$x_1,x_2,\ldots$ is again a completely continuous function of the new
variables after an orthogonal transformation.

If a bounded quadratic form $K$ is completely continuous, it is evident that
its eigenvalues do not accumulate in the finite plane; at the same time it
can be shown that an interval spectrum cannot be present at all.  From
\hyperref[thm:III]{Theorem III} we therefore obtain the following result.

\originalpage{201}
\begin{theorem}[V]\label{thm:V}
If a bounded form $K$ is completely continuous, it can always be brought by
an orthogonal substitution into the form
\begin{equation}
\tag{83}\label{eq:83}
  K(x)=k_1x_1^2+k_2x_2^2+\cdots;
\end{equation}
the quantities $k_1,k_2,\ldots$ are the reciprocal eigenvalues of $K$ and,
if infinitely many of them occur, have zero as their only condensation point.
\end{theorem}

Because of the many and important applications of this theorem, I give here
a very simple proof of it which is independent of the preceding theory.

\begin{proof}
First suppose that $K$ is a positive definite form.  Let
\[
  x_1=l_{11},\qquad x_2=l_{12},\qquad\ldots
\]
be values of the variables for which $K(x)$ attains the maximum $k_1$.
Clearly the sum of the squares of these values must equal $1$, since
otherwise we could increase the value of the form without violating the
condition $(x,x)\le1$.

Put
\[
  L_1(x)=l_{11}x_1+l_{12}x_2+\cdots
\]
and, now imposing on the variables the condition
\[
  L_1(x)=0,
\]
determine the relative maximum $k_2$ of $K(x)$.  Let it be attained for the
values
\[
  x_1=l_{21},\qquad x_2=l_{22},\qquad\ldots,
\]
whose sum of squares must again equal $1$.  Further put
\[
  L_2(x)=l_{21}x_1+l_{22}x_2+\cdots
\]
and, imposing on the variables the conditions
\[
  L_1(x)=0,\qquad L_2(x)=0,
\]
determine the relative maximum $k_3$ of $K(x)$.  Let it be attained for
\[
  x_1=l_{31},\qquad x_2=l_{32},\qquad\ldots.
\]
We then put
\[
  L_3(x)=l_{31}x_1+l_{32}x_2+\cdots.
\]
Continuing this procedure, we obtain a system of linear forms
$L_1,L_2,L_3,\ldots$ with the orthogonality properties
\[
  L_p(\mathord{\cdot})L_p(\mathord{\cdot})=1,
  \qquad
  L_p(\mathord{\cdot})L_q(\mathord{\cdot})=0\quad(p\ne q).
\]

\originalpage{202}
On the basis of the earlier considerations
(\hyperref[origpage:196]{p.~196}), complete these linear forms by a system of
linear forms
\[
  M_1(x),M_2(x),\ldots
\]
such that
\[
\begin{aligned}
  x'_1&=L_1(x),\\
  x'_2&=L_2(x),\\
      &\ \vdots\\
  y_1&=M_1(x),\\
  y_2&=M_2(x),\\
     &\ \vdots
\end{aligned}
\]
form an orthogonal substitution of the variables $x_1,x_2,\ldots$.  Denote
the form $K(x)$ transformed by this orthogonal substitution by $K(x'y)$.
The coefficient of ${x'_1}^{\,2}$ in $K(x'y)$ must plainly equal $k_1$.
On the other hand, no further terms containing $x'_1$ may occur in $K(x'y)$,
since the difference
\[
\begin{aligned}
  &K(x'y)-k_1\bigl({x'_1}^{\,2}+{x'_2}^{\,2}+\cdots
                   +y_1^2+y_2^2+\cdots\bigr)\\
  &\qquad=K(x)-k_1(x,x)
\end{aligned}
\]
is to be negative or zero for all values of the variables
$x'_1,x'_2,\ldots,y_1,y_2,\ldots$.  Since the same considerations apply to
$x'_2,x'_3,\ldots$, we have
\[
  K(x'y)=k_1{x'_1}^{\,2}+k_2{x'_2}^{\,2}+\cdots+R(y),
\]
where $R(y)$ denotes a quadratic form containing only the variables
$y_1,y_2,\ldots$.

Since $K$ is completely continuous, the same is true of the form
\[
  K(x'0)=k_1{x'_1}^{\,2}+k_2{x'_2}^{\,2}+\cdots;
\]
hence, if infinitely many of the quantities $k_1,k_2,\ldots$ occur, they must
converge to zero.  For otherwise there would be a sequence of values of
$K(x'0)$ converging to a value different from zero, while each of the
arguments $x'_1,x'_2,\ldots$ separately converged to zero.

Suppose now that there were a system of values
$y_1=m_1,y_2=m_2,\ldots$ for which $R(m)>0$.  We could then choose $q$ so
that $R(m)>k_q$.  For
\[
  x'_1=0,\quad x'_2=0,\quad\ldots,
  \qquad y_1=m_1,\quad y_2=m_2,\quad\ldots,
\]
we would also have $K>k_q$.  The equations
\[
  L_1(x)=0,\quad L_2(x)=0,\quad\ldots,
  \qquad M_1(x)=m_1,\quad M_2(x)=m_2,\quad\ldots
\]
would therefore determine a system of values of the variables
$x_1,x_2,\ldots$\originalpage{203} for which, in particular, the conditions
\[
  L_1(x)=0,\quad\ldots,\quad L_{q-1}(x)=0
\]
are satisfied and at the same time $K>k_q$.  This contradicts the manner in
which $k_q$ was determined; hence $R(y)$ cannot assume positive values.

Since
\[
  K(0y)=R(y),
\]
$R(y)$ certainly cannot assume negative values either, and consequently
$R(y)$ is identically zero; i.e.
\[
  K(x)=k_1L_1^2(x)+k_2L_2^2(x)+\cdots.
\]

If $K(x)$ is not assumed to be a definite form, the same reasoning leads to
the representation
\[
  K(x)=k_1L_1^2(x)+k_2L_2^2(x)+\cdots+R(y),
\]
where $R(y)$ cannot assume positive values.  Since $-R(y)$, as a positive
definite form, admits a representation of the same kind, we finally obtain
for $K$ also a representation by squares of orthogonal linear forms, and the
proof of \hyperref[thm:V]{Theorem V} is thereby complete.
\end{proof}

A sufficient criterion for the complete continuity of a form is provided by
the following theorem.

\begin{theorem}[VI]\label{thm:VI}
If, for a quadratic form $K$, one of the sums
\[
\begin{aligned}
  s_2&=\sum_{(p,q)}k_{pq}k_{qp}
      =\sum_{(p,q)}k_{pq}^{\,2},\\
  s_4&=\sum_{(p,q,r,s)}k_{pq}k_{qr}k_{rs}k_{sp},\\
     &\ \vdots
\end{aligned}
\]
remains finite, or if, for a definite quadratic form $K$, one of the sums
\[
\begin{aligned}
  s_1&=\sum_{(p)}k_{pp},\\
  s_3&=\sum_{(p,q,r)}k_{pq}k_{qr}k_{rp},\\
     &\ \vdots
\end{aligned}
\]
remains finite, then $K$ is a bounded, completely continuous form.
\end{theorem}

\begin{proof}
Indeed, if $K$ is a quadratic form whose coefficients have a finite sum of
squares, then, because
\[
\begin{aligned}
  \left\lvert
    x_1\frac{\partial K(x,y)}{\partial y_1}
    +x_2\frac{\partial K(x,y)}{\partial y_2}+\cdots
  \right\rvert^2
  &\le
  \left(\frac{\partial K(x,y)}{\partial y_1}\right)^2
  +\left(\frac{\partial K(x,y)}{\partial y_2}\right)^2+\cdots,\\[0.4em]
  \left(\frac{\partial K(x,y)}{\partial y_p}\right)^2
  &\le k_{p1}^2+k_{p2}^2+\cdots.
\end{aligned}
\]
\begin{samepage}
we necessarily have\originalpage{204}
\[
  \lvert K(x)\rvert
  \le\sqrt{\sum_{(p,q)}k_{pq}^{\,2}}.
\]
\end{samepage}
Applying this fact to the quadratic form
\[
  K(x)-K_n(x)=\sum_{(p,q)}^{(n)}k_{pq}x_p x_q,
\]
where on the right $(p,q)$ runs through all pairs of integral values except
those for which both $p\le n$ and $q\le n$, we find
\[
  \lvert K(x)-K_n(x)\rvert
  \le\sqrt{\sum_{(p,q)}^{(n)}k_{pq}^{\,2}}.
\]
Since
\[
  \mathop{\mathrm{L}}\limits_{n=\infty}
  \sqrt{\sum_{(p,q)}^{(n)}k_{pq}^{\,2}}=0,
\]
the required complete continuity of $K(x)$ follows.

If $K$ is a definite form, then
\[
  k_{pq}^2\le k_{pp}k_{qq},
\]
and hence
\[
  \sum_{(p,q)}k_{pq}^2\le\left(\sum_{(p)}k_{pp}\right)^2.
\]
Thus, if $\sum_{(p)}k_{pp}$ remains finite for a definite form, its
coefficients certainly also have a finite sum of squares, and, by the
preceding consideration, the form is again a completely continuous function
of the variables.

We now readily recognize, one after another, the following facts.

\noindent 1. If a bounded quadratic form $K$ is not completely continuous,
then it is also not completely continuous at the particular system of values
$0,0,\ldots$.  This assertion follows by applying the formula
\[
  K(x+a)=K(x)+2K(x,a)+K(a),
\]
with a system of values chosen for $a_1,a_2,\ldots$ at which $K$ is not
completely continuous, taking account of the fact that $K(x,a)$, as a bounded
linear form, is certainly completely continuous.

\noindent 2. If $K$ is a completely continuous quadratic form, then so is
the form $KK$; this follows immediately from \hyperref[thm:V]{Theorem V}.

\noindent 3. If $K$ is completely continuous, $K^*$ is any quadratic form,
and for all systems of values $x_1,x_2,\ldots$ the inequality
\[
  \lvert K^*(x)\rvert\le\lvert K(x)\rvert
\]

\originalpage{205}
holds, then $K^*$ is also completely continuous; for from this inequality
complete continuity for the system of values $0,0,\ldots$ follows.

4. If $K$ is completely continuous and $K^*$ is bounded, then the
convolution of $K$ and $K^*$ is completely continuous; for
\[
  \lvert K(x,\mathord{\cdot})K^*(x,\mathord{\cdot})\rvert
  \leq \sqrt{\bigl(KK(x)\bigr)\bigl(K^*K^*(x)\bigr)},
\]
this convolution is certainly completely continuous at the system of values
$0,0,\ldots$, since the form $KK$ is completely continuous by 2.

5. If the convolution $KK$ of a quadratic form $K$ is completely
continuous, then the form $K$ itself is also completely continuous; this
follows likewise, by virtue of 1, from the inequality
\[
  \lvert K(x)\rvert\leq\sqrt{KK(x)}.
\]

6. If one of the forms which have arisen by repeated convolution from the
bounded form $K$,
\[
  K^{(3)}=KKK,\qquad K^{(4)}=KKKK,\qquad
  K^{(5)}=KKKKK,\ldots.
\]
is completely continuous, then $K$ is also completely continuous. For if,
say, $K^{(f)}$ is completely continuous, then, by 4, the forms
$K^{(f+1)},K^{(f+2)},\ldots$ are also completely continuous. Among these we
choose a form for which the number of convolutions is a power of $2$, say
$K^{(2^g)}$; then, by applying 5 $g$ times, we conclude that $K$ is
completely continuous.

7. If the bounded quadratic form $K$ is definite, then the convolutions
$K^{(3)},K^{(5)},\ldots$ are likewise definite; for example, $K^{(3)}(x)$
arises from $K(x)$ when $x_p$ is replaced by
$\partial K(x,y)/\partial y_p$. Since, in general, $s_f$ is nothing other than
the invariant $K^{(f)}(\mathord{\cdot},\mathord{\cdot})$, namely the sum
of the coefficients of $x_p^2$ in $K^{(f)}$, it follows from 6 and 7, and
from the fact that the case $f=1$ has already been disposed of above, that
\hyperref[thm:VI]{Theorem VI} is valid in general.
\end{proof}

From \hyperref[thm:V]{Theorems V} and \hyperref[thm:VI]{VI} we obtain the
following fact:

\begin{theorem}[VII]\label{thm:VII}
A quadratic form $K$ satisfying one of the assumptions made in
\hyperref[thm:VI]{Theorem VI} can be transformed, by an orthogonal
transformation \eqref{eq:83}, into a sum of squares.
\end{theorem}

The counterpart to the case treated in \hyperref[thm:V]{Theorems V}--
\hyperref[thm:VII]{VII} is the assumption that $K$ has no point spectrum,
but only an interval spectrum. In order to consider only the simplest, yet
typical, case, we further make the assumptions
\originalpage{206}
that the interval spectrum $(s)$ consists of a finite number of intervals,
that the coefficients of the spectral form $\sigma(\mu,x)$ are continuously
differentiable with respect to $\mu$, and finally, when
\[
  \frac{d\sigma(\mu,x)}{d\mu}=\psi(\mu,x)
  =\sum_{p,q=1,2,\ldots}\psi_{pq}(\mu)x_px_q,
\]
that $\psi_{11}(\mu)$ is nowhere zero within $(s)$ and, on putting, for
brevity,
\[
  \psi_1(\mu)=\frac{\psi_{11}(\mu)}{\sqrt{\psi_{11}(\mu)}},
  \qquad
  \psi_2(\mu)=\frac{\psi_{12}(\mu)}{\sqrt{\psi_{11}(\mu)}},\ldots,
\]
the functions $\psi_1(\mu),\psi_2(\mu),\ldots$ are linearly independent in
the sense that, for an arbitrary function $u(\mu)$, there is no linear
relation among the integrals
\begin{equation}
  \int_{(s)}u(\mu)\psi_1(\mu)\,d\mu,\quad
  \int_{(s)}u(\mu)\psi_2(\mu)\,d\mu,\quad\ldots,
  \tag{84}\label{eq:84}
\end{equation}
whose coefficients are constants having a finite sum of squares.

If we insert our assumptions into relation \eqref{eq:82}, put
$\xi_1=1$, $\xi_2=0$, $\xi_3=0,\ldots$, and replace $u(\mu)$ by
$u(\mu)/\sqrt{\psi_{11}(\mu)}$, we obtain
\begin{equation}
  \int_{(s)}\bigl(u(\mu)\bigr)^2\,d\mu
  =\sum_{p=1,2,\ldots}
  \left\{\int_{(s)}u(\mu)\psi_p(\mu)\,d\mu\right\}^{\!2}.
  \tag{85}\label{eq:85}
\end{equation}
Consequently, for arbitrary functions $u(\mu)$ and $v(\mu)$,
\begin{equation}
  \int_{(s)}u(\mu)v(\mu)\,d\mu
  =\sum_{p=1}^{\infty}
  \int_{(s)}u(\mu)\psi_p(\mu)\,d\mu
  \int_{(s)}v(\mu)\psi_p(\mu)\,d\mu.
  \tag{86}\label{eq:86}
\end{equation}
If we take $v(\mu)=\psi_q(\mu)$ here, then
\begin{equation}
  \int_{(s)}u(\mu)\psi_q(\mu)\,d\mu
  =\sum_{p=1}^{\infty}
  \int_{(s)}\psi_p(\mu)\psi_q(\mu)\,d\mu
  \int_{(s)}\psi_p(\mu)u(\mu)\,d\mu.
  \tag{87}\label{eq:87}
\end{equation}
In view of the linear independence expressed in \eqref{eq:84}, relation
\eqref{eq:87} can hold identically for every $u(\mu)$ only when
\originalpage{207}
\begin{align}
  \int_{(s)}\bigl(\psi_p(\mu)\bigr)^2\,d\mu&=1,
  \nonumber\\
  \int_{(s)}\psi_p(\mu)\psi_q(\mu)\,d\mu&=0\qquad(p\ne q).
  \tag{88}\label{eq:88}
\end{align}
Now, from the positive definite character of the form $\psi(\mu,\xi)$,
\[
  \bigl(\psi_{11}\xi_1+\psi_{12}\xi_2+\cdots\bigr)^2
  \leq \psi(\mu,\xi)\psi_{11},
\]
or
\begin{equation}
  \bigl(\psi_1\xi_1+\psi_2\xi_2+\cdots\bigr)^2
  \leq \psi(\mu,\xi).
  \tag{89}\label{eq:89}
\end{equation}
On the other hand, by \eqref{eq:88},
\[
  \int_{(s)}
  \bigl(\psi_1\xi_1+\psi_2\xi_2+\cdots\bigr)^2\,d\mu
  =(\xi,\xi),
\]
and since
\[
  \int_{(s)}\psi(\mu,\xi)\,d\mu=(\xi,\xi),
\]
it follows that
\[
  \int_{(s)}
  \left\{
    \psi(\mu,\xi)
    -\bigl(\psi_1\xi_1+\psi_2\xi_2+\cdots\bigr)^2
  \right\}\,d\mu=0.
\]
But the expression under the integral sign is nowhere negative, by
\eqref{eq:89}; hence
\[
  \psi(\mu,\xi)
  =\bigl(\psi_1(\mu)\xi_1+\psi_2(\mu)\xi_2+\cdots\bigr)^2.
\]

Thus, under the assumptions made, the characteristic property of the
spectral form $\sigma(\mu,\xi)$ consists in the fact that its derivative with
respect to $\mu$ is the square of a linear form whose coefficients satisfy
the orthogonality relations \eqref{eq:85} and \eqref{eq:88}. Conversely,
since a form of this kind satisfies all the characteristic properties of a
spectral form, it is now easy to construct a quadratic form $K$ having any
prescribed interval spectrum $(s)$: we choose a complete orthogonal system of functions
$\psi_1(\mu),\psi_2(\mu),\ldots$ satisfying the relations \eqref{eq:85} and
\eqref{eq:88}, and put
\[
  K(\xi)=\int_{(s)}
  \frac{\bigl(\psi_1(\mu)\xi_1+\psi_2(\mu)\xi_2+\cdots\bigr)^2}{\mu}
  \,d\mu.
\]
The simplest example of this kind is the quadratic form
\begin{equation}
  K(x)=x_1x_2+x_2x_3+x_3x_4+\cdots;
  \tag{90}\label{eq:90}
\end{equation}
it has no point spectrum, while its interval spectrum
\originalpage{208}
consists of the two intervals from $\lambda=-\infty$ to $-1$ and from
$+1$ to $+\infty$. Here
\[
  \psi_p(\mu)=\sqrt{\frac{2}{\pi}}\,
  \frac{\cos t}{\sqrt{\sin t}}\sin pt,
\]
where $t$, lying between $0$ and $\pi$, is determined by
$\cos t=1/\mu$. We then have
\[
  \int_{\substack{(-\infty\cdots-1)\\(+1\cdots+\infty)}}
  \bigl(\psi_1(\mu)x_1+\psi_2(\mu)x_2+\cdots\bigr)^2\,d\mu
  =(x,x),
\]
and
\[
  K(x)=
  \int_{\substack{(-\infty\cdots-1)\\(+1\cdots+\infty)}}
  \frac{\bigl(\psi_1(\mu)x_1+\psi_2(\mu)x_2+\cdots\bigr)^2}{\mu}
  \,d\mu.
\]
In confirmation of \hyperref[thm:IV]{Theorem IV}, we note that the equations
\begin{align*}
  x_1-\frac{\lambda}{2}x_2&=0,\\
  x_2-\frac{\lambda}{2}(x_1+x_3)&=0,\\
  x_3-\frac{\lambda}{2}(x_2+x_4)&=0,\\
  &\ \vdots
\end{align*}
have, for no value of $\lambda$, solutions
$x_1,x_2,\ldots$ with finite sum of squares.

Another example with the same interval spectrum is the quadratic form
\begin{equation}
  \frac{2}{\sqrt{2^2-1}}x_1x_2
  +\frac{4}{\sqrt{4^2-1}}x_2x_3
  +\frac{6}{\sqrt{6^2-1}}x_3x_4+\cdots.
  \tag{91}\label{eq:91}
\end{equation}
For this form we find
\[
  \psi_p(\mu)=\sqrt{\frac{2p-1}{2}}\,\frac{1}{\mu}
  P^{(p-1)}\!\left(\frac{1}{\mu}\right),
\]
where $P^{(p-1)}(x)$ denotes the Legendre polynomial of degree $p-1$.
If we put
\[
  \chi_p(\mu)=\sqrt{\frac{2p-1}{2}}\,\frac{1}{\mu}
  Q^{(p-1)}\!\left(\frac{1}{\mu}\right),
\]
where $Q^{(p-1)}(x)$ denotes the associated spherical function of the
second kind, then
\originalpage{209}
the resolvent of the form \eqref{eq:91} has the form
\[
  K(\lambda,x)
  =\sum_{(p,\,q\geq p)}4\lambda\psi_p(\lambda)\chi_q(\lambda)x_px_q.
\]
The two forms \eqref{eq:90} and \eqref{eq:91} can be transformed into one
another by an orthogonal substitution, as is immediately seen from their
spectral forms. If we drop the assumption of linear independence made
above, the derivative of the spectral form is not the square of a single
linear form, but rather the sum of squares of linear forms.

\section*{XII. Simultaneous System of Quadratic Forms,\\
the Hermitian Form, the Skew-Symmetric Form,\\
and the Bilinear Form with Infinitely Many Variables}
\addcontentsline{toc}{section}{XII. Simultaneous System of Quadratic Forms, the Hermitian Form, the Skew-Symmetric Form, and the Bilinear Form with Infinitely Many Variables}
\label{sec:XII}

The methods and results developed in \hyperref[sec:XI]{Section XI} can be
extended, without any difficulty of principle, to more general forms with
infinitely many variables. We first consider the case of a simultaneous
system of two quadratic forms, one of which has definite character, while
the other is presented as an aggregate of positive and negative squares of
the variables. With the aid of my method of passage to the limit, starting
from forms with a finite number of variables, we can easily develop the
corresponding theory; we single out only the following result:

\begin{theorem}[VIII]\label{thm:VIII}
Let $K(x)$ be a positive definite, completely continuous, closed quadratic
form, and let
\[
  V(x)=v_1x_1^2+v_2x_2^2+\cdots,
\]
where each of the quantities $v_1,v_2,\ldots$ is either $+1$ or $-1$. Then
there exists an infinite sequence of quantities different from zero,
$\varkappa_1,\varkappa_2,\ldots$, having respectively the signs of
$v_1,v_2,\ldots$ and tending to zero. Their reciprocals are called the
eigenvalues of $K$ relative to $V$. There also exists an infinite sequence
of bounded linear forms $L_1(x),L_2(x),\ldots$, called the corresponding
eigenforms, such that the polar relations
\[
  L_p(\mathord{\cdot})V(\mathord{\cdot},\mathord{\cdot})
  L_p(\mathord{\cdot})=v_p,
  \qquad
  L_p(\mathord{\cdot})V(\mathord{\cdot},\mathord{\cdot})
  L_q(\mathord{\cdot})=0\quad(p\ne q)
\]
hold and
\originalpage{210}
\[
  K(x)=\lvert\varkappa_1\rvert\bigl(L_1(x)\bigr)^2
      +\lvert\varkappa_2\rvert\bigl(L_2(x)\bigr)^2+\cdots.
\]
\end{theorem}

This \hyperref[thm:VIII]{Theorem VIII} can also be proved, without passage
from a finite to an infinite number of variables, solely on the basis of
\hyperref[thm:V]{Theorem V}, with no new convergence considerations.

\begin{proof}
For this purpose we first suppose that the given form $K(x)$ is closed. We
then bring this form $K(x)$, by \hyperref[thm:V]{Theorem V}, through an
orthogonal transformation of the variables $x_1,x_2,\ldots$ into the form
of a sum of squares; we denote the new variables by $x'_1,x'_2,\ldots$ and
find
\[
  K(x)=k_1{x'_1}^{2}+k_2{x'_2}^{2}+\cdots,
\]
where $k_1,k_2,\ldots$ are all positive quantities converging to zero. We
further denote by $V'(x')$ the quadratic form in the variables
$x'_1,x'_2,\ldots$ which arises from $V(x)$ through that orthogonal
transformation, and finally by $V'(\sqrt{k}\,\xi)$ the quadratic form in the
variables $\xi_1,\xi_2,\ldots$ which arises from $V'(x')$ when
$\sqrt{k_1}\xi_1,\sqrt{k_2}\xi_2,\ldots$ are substituted in it for
$x'_1,x'_2,\ldots$.

We can now easily show that $V'(\sqrt{k}\,\xi)$ is a completely continuous
form in the variables $\xi_1,\xi_2,\ldots$. Indeed, as is readily seen,
$V'(x')$ can be represented as the difference of two elementary forms
$E_1(x')$ and $E_2(x')$; as such, they satisfy the inequalities
\[
  E_1(x')\leq(x',x'),\qquad E_2(x')\leq(x',x').
\]
Substituting $\sqrt{k_1}\xi_1,\sqrt{k_2}\xi_2,\ldots$ again for
$x'_1,x'_2,\ldots$, these inequalities become
\[
  E_1(\sqrt{k}\,\xi)\leq k_1\xi_1^2+k_2\xi_2^2+\cdots,
\]
\[
  E_2(\sqrt{k}\,\xi)\leq k_1\xi_1^2+k_2\xi_2^2+\cdots.
\]
If these elementary forms were not completely continuous in the variables
$\xi_1,\xi_2,\ldots$, there would have to exist systems of values
$\alpha_1^{(n)},\alpha_2^{(n)},\ldots$ for which
\[
  \lim_{n=\infty}\alpha_1^{(n)}=0,\qquad
  \lim_{n=\infty}\alpha_2^{(n)}=0,\ldots,
\]
while the elementary forms, which are positive definite, would have to
assume, for
\[
  \xi_1=\alpha_1^{(n)},\qquad \xi_2=\alpha_2^{(n)},\ldots,
\]
values remaining above a positive quantity independent of $n$;
\originalpage{211}
but this would contradict the foregoing inequalities, since $K(x)$ is
completely continuous. Since, accordingly,
$E_1(\sqrt{k}\,\xi)$ and $E_2(\sqrt{k}\,\xi)$ are completely continuous in
$\xi_1,\xi_2,\ldots$, so too is $V'(\sqrt{k}\,\xi)$. The complete continuity
of $E_1(\sqrt{k}\,\xi)$ and $E_2(\sqrt{k}\,\xi)$ in the variables
$\xi_1,\xi_2,\ldots$ also follows directly from
\hyperref[origpage:205]{4 on p.~205}.

We now transform, by \hyperref[thm:V]{Theorem V}, the variables
$\xi_1,\xi_2,\ldots$ orthogonally into the new variables
$\xi'_1,\xi'_2,\ldots$ in such a way that the form
$V'(\sqrt{k}\,\xi)$ assumes the form
\[
  V'(\sqrt{k}\,\xi)
  =\varkappa_1{\xi'_1}^{2}+\varkappa_2{\xi'_2}^{2}+\cdots,
\]
where $\varkappa_1,\varkappa_2,\ldots$ are certain real quantities
converging to zero. Let the linear forms in $x'_1,x'_2,\ldots$ which arise
from the forms $\xi'_p(\xi)$ when
$\sqrt{k_1}x'_1,\sqrt{k_2}x'_2,\ldots$ are substituted for
$\xi_1,\xi_2,\ldots$ be denoted by $\xi'_p(\sqrt{k}\,x')$. Since the former
forms constitute an orthogonal transformation, we then have
\begin{align*}
  K(x)
  &=\bigl(\sqrt{k_1}x'_1\bigr)^2
    +\bigl(\sqrt{k_2}x'_2\bigr)^2+\cdots\\
  &=\bigl(\xi'_1(\sqrt{k}\,x')\bigr)^2
    +\bigl(\xi'_2(\sqrt{k}\,x')\bigr)^2+\cdots;
\end{align*}
furthermore,
\begin{align*}
 &\xi'_p(\mathord{\cdot})
 V'(\sqrt{k}\,\xi,\sqrt{k}\,\mathord{\cdot}) \\
 &\quad=\frac{\partial\xi'_p(\xi)}{\partial\xi_1}
 \frac{\partial V'(\sqrt{k}\,\xi,\sqrt{k}\,\eta)}{\partial\eta_1}
 +\frac{\partial\xi'_p(\xi)}{\partial\xi_2}
 \frac{\partial V'(\sqrt{k}\,\xi,\sqrt{k}\,\eta)}{\partial\eta_2}
 +\cdots\\
 &\quad=\xi'_p(\sqrt{k}\,\mathord{\cdot})
 V'(\sqrt{k}\,\xi,\mathord{\cdot}),
\end{align*}
and consequently
\[
  \xi'_p(\sqrt{k}\,\mathord{\cdot})
  V'(\sqrt{k}\,\xi,\mathord{\cdot})
  =\varkappa_p\xi'_p(\xi).
\]
From this formula we infer in the same way
\[
  \xi'_p(\sqrt{k}\,\mathord{\cdot})V'(\mathord{\cdot},\mathord{\cdot})
  \xi'_q(\sqrt{k}\,\mathord{\cdot})
  =\varkappa_p\xi'_p(\mathord{\cdot})\xi'_q(\mathord{\cdot}).
\]
If we now return to the variables $x_1,x_2,\ldots$, and denote generally by
$\Lambda_p(x)$ the linear form in $x_1,x_2,\ldots$ which thereby arises from
$\xi'_p(\sqrt{k}\,x')$, we obtain
\[
  \Lambda_p(\mathord{\cdot})V(\mathord{\cdot},\mathord{\cdot})
  \Lambda_q(\mathord{\cdot})
  =\varkappa_p\xi'_p(\mathord{\cdot})\xi'_q(\mathord{\cdot}),
\]
and, by the orthogonality of the forms $\xi'_p(\xi)$, this is equal to
$\varkappa_p$ when $p=q$, and equal to zero when $p\ne q$; on the other hand,
\[
  K(x)=\Lambda_1^2(x)+\Lambda_2^2(x)+\cdots.
\]
If $\varkappa_p=0$, then for all $x_1,x_2,\ldots$ we should have
\[
  \Lambda_p(\mathord{\cdot})V(\mathord{\cdot},\mathord{\cdot})
  K(\mathord{\cdot},x)=0,
\]
which would contradict the closedness of $K$ and of $V$;
\originalpage{212}
consequently $\varkappa_1,\varkappa_2,\ldots$ are all different from zero.

We may now think of the terms
$\varkappa_1{\xi'_1}^{2},\varkappa_2{\xi'_2}^{2},\ldots$ in the above
representation of $V'(\sqrt{k}\,\xi)$ as arranged so that, in general,
$\varkappa_p$ has the sign of $v_p$. Indeed, such an arrangement of those
terms would fail to be possible only if either the number of negative
(respectively, positive) units in the sequence $v_1,v_2,\ldots$ were finite
and at the same time the number of negative (respectively, positive)
quantities occurring in the sequence
$\varkappa_1,\varkappa_2,\ldots$ were greater than the former number, or if
the number of negative (respectively, positive) quantities in the sequence
$\varkappa_1,\varkappa_2,\ldots$ were finite and at the same time the number
of negative (respectively, positive) units in the sequence
$v_1,v_2,\ldots$ were greater than the former number.

If we denote by
$x_1(\sqrt{k}\,\xi),x_2(\sqrt{k}\,\xi),\ldots$ the linear forms in
$\xi_1,\xi_2,\ldots$ which arise from $x_1(x'),x_2(x'),\ldots$ upon
substituting $\sqrt{k_1}\xi_1,\sqrt{k_2}\xi_2,\ldots$ for
$x'_1,x'_2,\ldots$, then, identically in $\xi_1,\xi_2,\ldots$,
\[
  v_1\bigl(x_1(\sqrt{k}\,\xi)\bigr)^2
  +v_2\bigl(x_2(\sqrt{k}\,\xi)\bigr)^2+\cdots
  =\varkappa_1\bigl(\xi'_1(\xi)\bigr)^2
  +\varkappa_2\bigl(\xi'_2(\xi)\bigr)^2+\cdots,
\]
because both sides of this equation represent $V'(\sqrt{k}\,\xi)$.

In accordance with the first case, suppose that $v_1,\ldots,v_e$ are
negative (respectively, positive), while
$v_{e+1},v_{e+2},\ldots$ are all positive (respectively, negative), and that
$\varkappa_1,\ldots,\varkappa_{e+1}$ are negative (respectively, positive).
Since the forms $\xi'_1(\xi),\xi'_2(\xi),\ldots$ determine an orthogonal
substitution, that is, form a \emph{complete orthogonal system of linear
forms}--as we shall say--every bounded linear form in
$\xi_1,\xi_2,\ldots$ can be represented as a linear combination of the
forms $\xi'_1(\xi),\xi'_2(\xi),\ldots$. In particular, we put
\begin{align*}
  x_1(\sqrt{k}\,\xi)
    &=a_{11}\xi'_1(\xi)+a_{12}\xi'_2(\xi)+\cdots,\\
  &\ \vdots\\
  x_e(\sqrt{k}\,\xi)
    &=a_{e1}\xi'_1(\xi)+a_{e2}\xi'_2(\xi)+\cdots,
\end{align*}
where $a_{11},a_{12},\ldots,a_{e1},a_{e2},\ldots$ denote certain
coefficients. We then choose quantities
$\alpha_1,\ldots,\alpha_{e+1}$, not all zero, which satisfy the $e$ equations
\begin{align*}
  a_{11}\alpha_1+\cdots+a_{1,e+1}\alpha_{e+1}&=0,\\
  &\ \vdots\\
  a_{e1}\alpha_1+\cdots+a_{e,e+1}\alpha_{e+1}&=0,
\end{align*}
and form the equations
\originalpage{213}
\begin{align*}
  \xi'_1(\xi)&=\alpha_1,\\
  &\ \vdots\\
  \xi'_{e+1}(\xi)&=\alpha_{e+1},\\
  \xi'_{e+2}(\xi)&=0,\\
  \xi'_{e+3}(\xi)&=0,\\
  &\ \vdots
\end{align*}
The values of $\xi_1,\xi_2,\ldots$ obtained by solving these equations would
give a contradiction, since, when inserted into the identity set up above,
\[
  v_1\bigl(x_1(\sqrt{k}\,\xi)\bigr)^2
  +v_2\bigl(x_2(\sqrt{k}\,\xi)\bigr)^2+\cdots
  =\varkappa_1\bigl(\xi'_1(\xi)\bigr)^2
  +\varkappa_2\bigl(\xi'_2(\xi)\bigr)^2+\cdots,
\]
they would give the left-hand side a nonnegative (respectively,
nonpositive) value, but the right-hand side certainly a negative
(respectively, positive) value.

To treat the second of the cases named above, we start from the equation,
identical in $x'_1,x'_2,\ldots$,
\[
  v_1\bigl(x_1(x')\bigr)^2+v_2\bigl(x_2(x')\bigr)^2+\cdots
  =\varkappa_1\left(\xi'_1\!\left(\frac{x'}{\sqrt{k}}\right)\right)^2
  +\varkappa_2\left(\xi'_2\!\left(\frac{x'}{\sqrt{k}}\right)\right)^2
  +\cdots,
\]
where
$\xi'_1(x'/\sqrt{k}),\xi'_2(x'/\sqrt{k}),\ldots$ denote the linear forms
in $x'_1,x'_2,\ldots$ which arise from
$\xi'_1(\xi),\xi'_2(\xi),\ldots$ upon substituting
$x'_1/\sqrt{k_1},x'_2/\sqrt{k_2},\ldots$ for
$\xi_1,\xi_2,\ldots$. Since, however, the linear forms
$\xi'_1(x'/\sqrt{k}),\xi'_2(x'/\sqrt{k}),\ldots$ do not necessarily become
bounded linear forms in $x'_1,x'_2,\ldots$, the above identity has meaning
only as an identity of truncations, and the procedure used in the first case
requires the following modification.

Suppose that $\varkappa_1,\ldots,\varkappa_e$ are negative (respectively,
positive). Suppose further that $\varkappa_{e+1},\varkappa_{e+2},\ldots$ are
all positive (respectively, negative), and that $v_1,\ldots,v_{e+1}$ are
negative (respectively, positive). We then put
\begin{align*}
  \xi'_1\!\left(\frac{x'}{\sqrt{k}}\right)
    &=\alpha_{11}x'_1+\alpha_{12}x'_2+\cdots,\\
  &\ \vdots\\
  \xi'_e\!\left(\frac{x'}{\sqrt{k}}\right)
    &=\alpha_{e1}x'_1+\alpha_{e2}x'_2+\cdots.
\end{align*}
Further, we think of the equations
\[
  x_1(x')=a_1,\ldots,x_{e+1}(x')=a_{e+1},\quad
  x_{e+2}(x')=0,\quad x_{e+3}(x')=0,\ldots
\]
as solved for $x'_1,x'_2,\ldots$, and represent the solutions as functions
\originalpage{214}
of $a_1,\ldots,a_{e+1}$ as follows:
\begin{align*}
  x'_1&=o_{11}a_1+\cdots+o_{1,e+1}a_{e+1},\\
  x'_2&=o_{12}a_1+\cdots+o_{2,e+1}a_{e+1},\\
  &\ \vdots
\end{align*}
Finally, for each $n$ we determine $e+1$ quantities
$a_1^{(n)},\ldots,a_{e+1}^{(n)}$ such that, after these values of
$x'_1,x'_2,\ldots$ have been inserted, the $e+1$ equations
\begin{align*}
  \alpha_{11}x'_1+\cdots+\alpha_{1n}x'_n&=0,\\
  &\ \vdots\\
  \alpha_{e1}x'_1+\cdots+\alpha_{en}x'_n&=0,\\
  a_1^2+\cdots+a_{e+1}^2&=1
\end{align*}
are satisfied for
\[
  a_1=a_1^{(n)},\ldots,a_{e+1}=a_{e+1}^{(n)}.
\]
We now select integers $n=n_h$ such that the limits of
$a_1^{(n_h)},\ldots,a_{e+1}^{(n_h)}$ exist as $h\to\infty$, and insert the
quantities $x'_1,\ldots,x'_{n_h}$ determined by these values
$a_1^{(n_h)},\ldots,a_{e+1}^{(n_h)}$ into the $n_h$th truncation of the above
identity
\[
  v_1\bigl(x_1(x')\bigr)^2+v_2\bigl(x_2(x')\bigr)^2+\cdots
  =\varkappa_1\left(\xi'_1\!\left(\frac{x'}{\sqrt{k}}\right)\right)^2
  +\varkappa_2\left(\xi'_2\!\left(\frac{x'}{\sqrt{k}}\right)\right)^2
  +\cdots.
\]
We then see that the left-hand side of this identity, since it represents a
bounded form in the variables $x'_1,x'_2,\ldots$ and, as such, is continuous
in these variables by \hyperref[origpage:178]{p.~178}, tends to the value
$-1$ as $h\to\infty$, whereas the right-hand side is always $\geq0$.

Both cases have thus been recognized as impossible, and we may therefore
assume from the outset, in general, that $\varkappa_p$ has the same sign as
$v_p$. If we now put
\[
  L_p(x)=\frac{\Lambda_p(x)}{\sqrt{\lvert\varkappa_p\rvert}},
\]
the linear forms $L_1(x),L_2(x),\ldots$ have the properties required in
\hyperref[thm:VIII]{Theorem VIII}.

The same mode of inference permits the treatment of a form $K$ which is not
closed. To see this, we again bring the form $K$, by
\hyperref[thm:V]{Theorem V}, through an orthogonal transformation of the
variables $x_1,
\originalpage{215}
x_2,\ldots$ into the form of a sum of squares. We put
\[
  K(x)=k_1{x'_1}^{2}+k_2{x'_2}^{2}+\cdots,
\]
where $k_1,k_2,\ldots$ are quantities of which some are positive and some
vanish. We again denote by $V'(x')$ the quadratic form in the variables
$x'_1,x'_2,\ldots$ arising from $V(x)$ through that orthogonal
transformation, and finally by $V'(\sqrt{k}\,\xi)$ the quadratic form in the
variables $\xi_1,\xi_2,\ldots$ which arises from $V'(x')$ when
$\sqrt{k_1}\xi_1,\sqrt{k_2}\xi_2,\ldots$ are substituted for
$x'_1,x'_2,\ldots$. Since $V'(\sqrt{k}\,\xi)$ is a completely continuous
form in $\xi_1,\xi_2,\ldots$, we may, by \hyperref[thm:V]{Theorem V},
transform the variables $\xi_1,\xi_2,\ldots$ orthogonally into the new
variables $\xi'_1,\xi'_2,\ldots$ in such a way that
\[
  V'(\sqrt{k}\,\xi)
  =\varkappa_1{\xi'_1}^{2}+\varkappa_2{\xi'_2}^{2}+\cdots,
\]
where $\varkappa_1,\varkappa_2,\ldots$ are certain quantities, each positive,
negative, or zero, which, if infinitely many occur, converge to zero. Finally,
forming as before the expressions $\xi'_p(\sqrt{k}\,x')$, and denoting
generally by $\Lambda_p(x)$ the linear form
which arises from $\xi'_p(\sqrt{k}\,x')$ when the original variables
$x_1,x_2,\ldots$ are reintroduced in place of the variables
$x'_1,x'_2,\ldots$, we obtain, as before,
\[
  K(x)=\Lambda_1^2(x)+\Lambda_2^2(x)+\cdots,
\]
\[
  \Lambda_p(\mathord{\cdot})V(\mathord{\cdot},\mathord{\cdot})
  \Lambda_q(\mathord{\cdot})
  =\begin{cases}
     0, & p\ne q,\\
     \varkappa_p, & p=q.
   \end{cases}
\]
\end{proof}

We state this result, which supplements \hyperref[thm:VIII]{Theorem VIII},
as follows:

\begin{theorem}[VIII*]\label{thm:VIIIstar}
Let a positive definite, completely continuous quadratic form $K(x)$, and
also a quadratic form of the shape
\[
  V(x)=v_1x_1^2+v_2x_2^2+\cdots,
\]
be given, where $v_1,v_2,\ldots$ have prescribed values $+1$ or $-1$. Then
there always exists a sequence of quantities
$\varkappa_1,\varkappa_2,\ldots$, each positive, negative, or zero and tending
to zero if infinitely many occur, together with associated bounded linear
forms $\Lambda_1(x),\Lambda_2(x),\ldots$, such that the
polar relations
\[
  \Lambda_p(\mathord{\cdot})V(\mathord{\cdot},\mathord{\cdot})
  \Lambda_p(\mathord{\cdot})=\varkappa_p,
\]
\[
  \Lambda_p(\mathord{\cdot})V(\mathord{\cdot},\mathord{\cdot})
  \Lambda_q(\mathord{\cdot})=0\qquad(p\ne q)
\]
hold, and such that, moreover, the given quadratic form permits the
representation
\[
  K(x)=\Lambda_1^2(x)+\Lambda_2^2(x)+\cdots.
\]
\end{theorem}

\begin{center}
  \rule{0.22\linewidth}{0.4pt}
\end{center}

\originalpage{216}
By a \emph{Hermitian form} in the infinitely many variables
$x_1,x_2,\ldots,y_1,y_2,\ldots$ we understand a bilinear form in these
variables of the shape
\[
  H(x,y)=\sum_{p,q}h_{pq}x_py_q,
\]
whose coefficients $h_{pq}$ are complex quantities satisfying the condition
\[
  h_{pq}=k_{pq}+is_{pq}=\overline{h}_{qp}=k_{qp}-is_{qp}.
\]
If both the real and the imaginary parts of $H(x,y)$ are completely
continuous functions of the real variables
$x_1,x_2,\ldots,y_1,y_2,\ldots$, then real values
$\lambda_1,\lambda_2,\ldots$ having no finite condensation point--the
eigenvalues of $H$--and corresponding linear forms with complex coefficients
$L_1(x),L_2(x),\ldots$--the eigenforms of $H$--can be found such that
\[
  (x,y)=L_1(x)\overline{L}_1(y)+L_2(x)\overline{L}_2(y)+\cdots,
\]
\[
  H(x,y)=\frac{L_1(x)\overline{L}_1(y)}{\lambda_1}
         +\frac{L_2(x)\overline{L}_2(y)}{\lambda_2}+\cdots,
\]
and such that the orthogonality properties
\[
  L_p(\mathord{\cdot})\overline{L}_p(\mathord{\cdot})=1,
  \qquad
  L_p(\mathord{\cdot})\overline{L}_q(\mathord{\cdot})=0
  \quad(p\ne q)
\]
are satisfied; the horizontal bars indicate the interchange of $i$ with
$-i$. The proof of this fact can be conducted analogously to the proof,
given below, of the more special \hyperref[thm:IX]{Theorem IX}.

\begin{center}
  \rule{0.22\linewidth}{0.4pt}
\end{center}

If we assume the coefficients of the Hermitian form to be purely imaginary
and then suppress the factor $i$, the skew-symmetric form results. By a
\emph{skew-symmetric form} we therefore understand a bilinear form in the
variables $x_1,x_2,\ldots,y_1,y_2,\ldots$ of the shape
\[
  S(x,y)=\sum_{p,q}s_{pq}x_py_q,
\]
whose coefficients are real quantities satisfying the conditions
\[
  s_{pq}=-s_{qp},\qquad s_{pp}=0.
\]
The fact stated above for a Hermitian form is expressed, for the special case
of the skew-symmetric form, as follows:

\begin{theorem}[IX]\label{thm:IX}
If the skew-symmetric form $S(x,y)$ is completely continuous, then there is
an orthogonal transformation of the variables
\[
  x_1,x_2,x_3,x_4,\ldots
\]
\originalpage{217}
into the new variables
\[
  \xi_1,\xi'_1,\xi_2,\xi'_2,\ldots,
\]
such that, when the variables $y_1,y_2,y_3,y_4,\ldots$ are simultaneously
transformed by the same orthogonal transformation into the variables
$\eta_1,\eta'_1,\eta_2,\eta'_2,\ldots$, the form $S$ assumes the shape
\[
  S(x,y)=k_1(\xi_1\eta'_1-\xi'_1\eta_1)
        +k_2(\xi_2\eta'_2-\xi'_2\eta_2)+\cdots;
\]
here $k_1,k_2,\ldots$ are quantities which, if infinitely many occur, have
zero as their sole condensation value.
\end{theorem}

\begin{proof}
We regard $2S(x,y)$ as a quadratic form in the infinitely many variables
$x_1,y_1,x_2,y_2,\ldots$. From \hyperref[thm:V]{Theorem V} we then infer the
existence of quantities $k_1,k_2,\ldots$ and corresponding linear forms
$L_1(x,y),L_2(x,y),\ldots$ in those variables, such that
\[
  2S(x,y)=k_1\bigl(L_1(x,y)\bigr)^2
          +k_2\bigl(L_2(x,y)\bigr)^2+\cdots
\]
and
\[
  (x,x)+(y,y)=\bigl(L_1(x,y)\bigr)^2
              +\bigl(L_2(x,y)\bigr)^2+\cdots.
\]
In view of the properties of the skew-symmetric form
\[
  S(x,y)=-S(y,x),
  \qquad
  S(x,y)=-S(-x,y),
\]
it follows easily that, in the above representation of $2S(x,y)$, for every
$k_p,L_p$ there must also occur the eigenvalues and corresponding linear
forms
\begin{align*}
  k_{p'}&=-k_p,
  &L_{p'}(x,y)&=L_p(y,x),\\
  k_{p''}&=-k_p,
  &L_{p''}(x,y)&=L_p(-x,y),\\
  k_{p'''}&=-k_{p'}=k_p,
  &L_{p'''}(x,y)&=L_{p'}(-x,y)=L_p(y,-x).
\end{align*}
Their union contributes to the representation of $2S(x,y)$ the terms
\[
  k_p\left\{
    \bigl(L_p(x,y)\bigr)^2-\bigl(L_p(y,x)\bigr)^2
    -\bigl(L_p(-x,y)\bigr)^2+\bigl(L_p(y,-x)\bigr)^2
  \right\},
\]
and to the representation of $(x,x)+(y,y)$ the terms
\[
  \bigl(L_p(x,y)\bigr)^2+\bigl(L_p(y,x)\bigr)^2
  +\bigl(L_p(-x,y)\bigr)^2+\bigl(L_p(y,-x)\bigr)^2.
\]
If we now put
\[
  L_p(x,y)=\frac{1}{\sqrt2}\bigl(O_p(x)+O'_p(y)\bigr),
\]
where $O_p(x)$ is a linear form in $x_1,x_2,\ldots$ and $O'_p(y)$ a linear
form in $y_1,y_2,\ldots$, the above representations become
\[
  S(x,y)=\sum_p k_p\bigl(O_p(x)O'_p(y)-O'_p(x)O_p(y)\bigr),
\]
\[
  (x,x)+(y,y)=\sum_p\left\{
    \bigl(O_p(x)\bigr)^2+\bigl(O'_p(x)\bigr)^2
    +\bigl(O_p(y)\bigr)^2+\bigl(O'_p(y)\bigr)^2
  \right\},
\]
\originalpage{218}
and, since $L_1(x,y),L_2(x,y),\ldots$ are mutually orthogonal, the
orthogonality of the forms $O_1,O_2,\ldots$ follows readily as well. Thus
\[
  \xi_p=O_p(x),\qquad \xi'_p=O'_p(x)
\]
determine an orthogonal transformation of the required kind.
\end{proof}

From this representation it follows, by a simple consideration of the kind
that will be made similarly below on \hyperref[origpage:224]{p.~224}, that
the inhomogeneous equations arising from
$(x,y)-\lambda S(x,y)$ are uniquely solvable except when
\[
  \lambda=\frac{i}{k_1},\frac{i}{k_2},\ldots.
\]
For these purely imaginary eigenvalues of the form $S(x,y)$, the homogeneous
equations have a solution not identically zero, and the number of linearly
independent solutions is always finite.

\begin{center}
  \rule{0.22\linewidth}{0.4pt}
\end{center}

As regards, finally, the theory of the bilinear form, we first perceive the
following facts without difficulty. If the bilinear form $A(x,y)$ is a
completely continuous function of the infinitely many variables
$x_1,x_2,\ldots,y_1,y_2,\ldots$, then, if $A_n$ denotes the $n$th truncation
of the bilinear form $A$, for every system of values of the infinitely many
variables
\[
  \lim_{n=\infty}
  A_n(\mathord{\cdot},x)A_n(\mathord{\cdot},x)
  =A(\mathord{\cdot},x)A(\mathord{\cdot},x),
\]
and indeed in the sense of uniform convergence; that is,
\begin{equation}
  \left\lvert
  A(\mathord{\cdot},x)A(\mathord{\cdot},x)
  -A_n(\mathord{\cdot},x)A_n(\mathord{\cdot},x)
  \right\rvert\leq\varepsilon_n,
  \tag{92}\label{eq:92}
\end{equation}
where $\varepsilon_n$ are certain quantities independent of the variables
$x_1,x_2,\ldots$ and tending to zero as $n$ increases indefinitely. It
follows that the quadratic form
$A(\mathord{\cdot},x)A(\mathord{\cdot},x)$ is always completely continuous
when the bilinear form $A(x,y)$ is completely continuous.

A bilinear form $A(x,y)$ is always completely continuous if the quadratic
form $A(\mathord{\cdot},x)A(\mathord{\cdot},x)$ is completely continuous;
thus this is certainly the case, for example, when the sum of the squares of
the coefficients of $A$ remains finite. Indeed, if we regard $A(x,y)$ as a
quadratic form in the variables $x_1,x_2,\ldots,y_1,y_2,\ldots$, then from
the inequality
\[
  \lvert A(x,y)\rvert
  \leq\sqrt{A(\mathord{\cdot},x)A(\mathord{\cdot},x)},
\]
it follows, by the reasoning employed in 4 and 5 on
\hyperref[origpage:205]{p.~205}, that $A(x,y)$ is completely continuous.

I should like to close this communication with the development of a theorem
which--as I shall show in the following communication--
\originalpage{219}
can be used in the simplest manner for solving integral equations of the
second kind with a nonsymmetric kernel. It reads:

\begin{theorem}[X]\label{thm:X}
If
\[
  A(x,y)=\sum_{p,q}a_{pq}x_py_q
\]
is a completely continuous bilinear form in the infinitely many variables
$x_1,x_2,\ldots$ and $y_1,y_2,\ldots$, then certainly either the infinitely many
equations
\begin{equation}
  \begin{aligned}
    (1+a_{11})x_1+a_{12}x_2+\cdots&=a_1,\\
    a_{21}x_1+(1+a_{22})x_2+\cdots&=a_2,\\
    &\ \vdots
  \end{aligned}
  \tag{93}\label{eq:93}
\end{equation}
have, for all possible quantities $a_1,a_2,\ldots$ with convergent sum of
squares, a uniquely determined solution $x_1,x_2,\ldots$ with convergent sum
of squares--or the corresponding homogeneous equations
\begin{equation}
  \begin{aligned}
    (1+a_{11})x_1+a_{12}x_2+\cdots&=0,\\
    a_{21}x_1+(1+a_{22})x_2+\cdots&=0,\\
    &\ \vdots
  \end{aligned}
  \tag{94}\label{eq:94}
\end{equation}
admit a solution $x_1,x_2,\ldots$ having sum of squares $1$.
\end{theorem}

\begin{proof}
We first consider any system of $n$ linear equations in $n$ unknowns, with
nonvanishing determinant, of the form
\begin{align*}
  b_{11}x_1+\cdots+b_{1n}x_n&=b_1,\\
  &\ \vdots\\
  b_{n1}x_1+\cdots+b_{nn}x_n&=b_n.
\end{align*}
If $\beta_1,\ldots,\beta_n$ denote the solutions of these equations, then
certainly
\begin{align*}
 &(b_{11}\beta_1+\cdots+b_{1n}\beta_n)^2+\cdots
 +(b_{n1}\beta_1+\cdots+b_{nn}\beta_n)^2\\
 &\quad=b_1(b_{11}\beta_1+\cdots+b_{1n}\beta_n)+\cdots
 +b_n(b_{n1}\beta_1+\cdots+b_{nn}\beta_n),
\end{align*}
and consequently
\begin{multline*}
 (b_{11}\beta_1+\cdots+b_{1n}\beta_n)^2+\cdots
 +(b_{n1}\beta_1+\cdots+b_{nn}\beta_n)^2\\
 \leq\sqrt{(\beta_1^2+\cdots+\beta_n^2)
 \left\{
 (b_{11}b_1+\cdots+b_{n1}b_n)^2+\cdots
 +(b_{1n}b_1+\cdots+b_{nn}b_n)^2
 \right\}}.
\end{multline*}
Let $m$ now be the minimum of the quadratic form
\[
  (b_{11}x_1+\cdots+b_{1n}x_n)^2+\cdots
  +(b_{n1}x_1+\cdots+b_{nn}x_n)^2
\]
subject to the subsidiary condition
\begin{equation}
  x_1^2+\cdots+x_n^2=1,
  \tag{95}\label{eq:95}
\end{equation}
\originalpage{220}
and let $M$ be the maximum of the quadratic form
\[
  (b_{11}x_1+\cdots+b_{n1}x_n)^2+\cdots
  +(b_{1n}x_1+\cdots+b_{nn}x_n)^2
\]
under the same subsidiary condition. The preceding inequality then implies
the inequality
\begin{equation}
  \beta_1^2+\cdots+\beta_n^2
  \leq(b_1^2+\cdots+b_n^2)\frac{M}{m^2}.
  \tag{96}\label{eq:96}
\end{equation}

We apply this result to the equations
\begin{equation}
  \begin{aligned}
    (1+a_{11})x_1+\cdots+a_{1n}x_n&=a_1,\\
    &\ \vdots\\
    a_{n1}x_1+\cdots+(1+a_{nn})x_n&=a_n,
  \end{aligned}
  \tag{97}\label{eq:97}
\end{equation}
and, for this purpose, denote by $m_n$ the minimum of the quadratic form
\[
  \bigl((1+a_{11})x_1+\cdots+a_{1n}x_n\bigr)^2+\cdots
  +\bigl(a_{n1}x_1+\cdots+(1+a_{nn})x_n\bigr)^2
\]
subject to the subsidiary condition \eqref{eq:95}, and by $M_n$ the maximum
of the quadratic form
\[
  \bigl((1+a_{11})x_1+\cdots+a_{n1}x_n\bigr)^2+\cdots
  +\bigl(a_{1n}x_1+\cdots+(1+a_{nn})x_n\bigr)^2
\]
under the same subsidiary condition.

Because of the assumed complete continuity of the bilinear form $A(x,y)$,
the form in the infinitely many variables $x_1,x_2,\ldots$
\[
  \bigl((1+a_{11})x_1+a_{21}x_2+\cdots\bigr)^2
  +\bigl(a_{12}x_1+(1+a_{22})x_2+\cdots\bigr)^2+\cdots
\]
is certainly a bounded form, and from this we see that the maxima $M_n$ also
remain below a finite quantity $M$ independent of $n$.

As regards the minima $m_n$, two cases must be distinguished.

First, suppose that $m_n>m>0$ for infinitely many indices $n$. Then, for each
such $n$, equations
\eqref{eq:97} are soluble, and it follows from \eqref{eq:96} that the sum of
squares of their solutions lies below the finite quantity, independent of
$n$,
\begin{equation}
  (a,a)\frac{M}{m^2}.
  \tag{98}\label{eq:98}
\end{equation}
If, for such $n$, we denote the solutions of \eqref{eq:97} by
\[
  \alpha_1^{(n)},\ldots,\alpha_n^{(n)},
\]
then, by an often-used procedure, we can select from those $n$ integers
$n_1,n_2,\ldots$ such
\originalpage{221}
that the limits
\[
  \alpha_1=\lim_{h=\infty}\alpha_1^{(n_h)},\qquad
  \alpha_2=\lim_{h=\infty}\alpha_2^{(n_h)},\ldots
\]
exist. The quantities $\alpha_1,\alpha_2,\ldots$ have a sum of squares which
likewise lies below the bound \eqref{eq:98}, and, because of the complete
continuity of the linear forms
\[
  a_{p1}x_1+a_{p2}x_2+\cdots,
\]
they must certainly satisfy every equation of the given system
\eqref{eq:93}.

Second, let the minima $m_1,m_2,\ldots$ converge to zero; let the values of
the variables, with sum of squares $1$, at which these minima occur be
\[
  \mu_1^{(n)},\ldots,\mu_n^{(n)}.
\]
\begin{samepage}
For these values,
\begin{align*}
 &\bigl((1+a_{11})x_1+\cdots+a_{1n}x_n\bigr)^2+\cdots
 +\bigl(a_{n1}x_1+\cdots+(1+a_{nn})x_n\bigr)^2\\
 &\qquad=A_n(\mathord{\cdot},x)A_n(\mathord{\cdot},x)
 +2A_n(x,x)+(x,x)_n,
\end{align*}
\end{samepage}
and consequently
\begin{equation}
  A_n(\mathord{\cdot},\mu^{(n)})A_n(\mathord{\cdot},\mu^{(n)})
  +2A_n(\mu^{(n)},\mu^{(n)})+1=m_n.
  \tag{99}\label{eq:99}
\end{equation}
We again select a sequence of integers $n_1,n_2,\ldots$ such that the limits
\[
  \mu_1=\lim_{h=\infty}\mu_1^{(n_h)},\qquad
  \mu_2=\lim_{h=\infty}\mu_2^{(n_h)},\ldots
\]
exist. The quantities $\mu_1,\mu_2,\ldots$ then satisfy the condition
\begin{equation}
  (\mu,\mu)\leq1.
  \tag{100}\label{eq:100}
\end{equation}
In view of \eqref{eq:92}, and because of the complete continuity of the
quadratic form $A(x,x)$, it follows from \eqref{eq:99}, upon putting $n_h$
in place of $n$ and passing to the limit $h\to\infty$, that
\begin{equation}
  A(\mathord{\cdot},\mu)A(\mathord{\cdot},\mu)
  +2A(\mu,\mu)+1=0.
  \tag{101}\label{eq:101}
\end{equation}
We now consider the quadratic form
\begin{equation}
  \begin{aligned}
   &\bigl((1+a_{11})x_1+a_{12}x_2+\cdots\bigr)^2
   +\bigl(a_{21}x_1+(1+a_{22})x_2+\cdots\bigr)^2+\cdots\\
   &\qquad=A(\mathord{\cdot},x)A(\mathord{\cdot},x)
   +2A(x,x)+(x,x);
  \end{aligned}
  \tag{102}\label{eq:102}
\end{equation}
since it is positive definite, it follows in particular that
\begin{equation}
  A(\mathord{\cdot},\mu)A(\mathord{\cdot},\mu)
  +2A(\mu,\mu)+(\mu,\mu)\geq0.
  \tag{103}\label{eq:103}
\end{equation}
From this, in view of \eqref{eq:101}, we obtain
\[
  (\mu,\mu)\geq1;
\]
\originalpage{222}
hence, by \eqref{eq:100},
\[
  (\mu,\mu)=1.
\]
We now see from \eqref{eq:101} that also
\[
  A(\mathord{\cdot},\mu)A(\mathord{\cdot},\mu)
  +2A(\mu,\mu)+(\mu,\mu)=0;
\]
that is, in view of \eqref{eq:102}, the quantities
$\mu_1,\mu_2,\ldots$ satisfy the homogeneous equations \eqref{eq:94}. It has
therefore been shown that at least one of the cases distinguished in
\hyperref[thm:X]{Theorem X} always occurs.

If the homogeneous equations \eqref{eq:94} possess a solution with sum of
squares $1$, then the inhomogeneous equations obtained by transposition,
\begin{equation}
  \begin{aligned}
    (1+a_{11})x_1+a_{21}x_2+\cdots&=a_1,\\
    a_{12}x_1+(1+a_{22})x_2+\cdots&=a_2,\\
    &\ \vdots
  \end{aligned}
  \tag{104}\label{eq:104}
\end{equation}
certainly cannot, for all $a_1,a_2,\ldots$ with finite sum of squares,
possess a solution of finite sum of squares, since a linear identity exists
among their left-hand sides. By what has just been proved, therefore, the
transposed homogeneous equations
\begin{equation}
  \begin{aligned}
    (1+a_{11})x_1+a_{21}x_2+\cdots&=0,\\
    a_{12}x_1+(1+a_{22})x_2+\cdots&=0,\\
    &\ \vdots
  \end{aligned}
  \tag{105}\label{eq:105}
\end{equation}
must then admit a solution with sum of squares $1$. Thus the inhomogeneous
equations \eqref{eq:93} certainly cannot possess a solution of finite sum of
squares for all $a_1,a_2,\ldots$. The two cases in
\hyperref[thm:X]{Theorem X} are therefore indeed mutually exclusive, and the
solution in the first case is unique. The proof of our theorem is thereby
fully completed.
\end{proof}

To determine the multiplicity of the solutions of the homogeneous equations
\eqref{eq:94}, we need only apply to the form \eqref{eq:102} the theory of
the orthogonal transformation of quadratic forms developed in
\hyperref[sec:XI]{Section XI}. Since
$A(\mathord{\cdot},x)A(\mathord{\cdot},x)$ and $A(x,x)$ are completely
continuous quadratic forms, so too is the form
\[
  A(\mathord{\cdot},x)A(\mathord{\cdot},x)+2A(x,x);
\]
the latter therefore possesses the value $-1$ at most as an eigenvalue of
finite multiplicity. Consequently, the quadratic form \eqref{eq:102}
possesses the value $\infty$ only as an eigenvalue of finite multiplicity;
that is, there is an orthogonal transformation of the variables
\originalpage{223}
$x_1,x_2,\ldots$ into $x'_1,x'_2,\ldots$ such that the quadratic form
\eqref{eq:102} assumes the shape
\[
  k_e{x'_e}^{2}+k_{e+1}{x'_{e+1}}^{2}+\cdots,
\]
where $k_e,k_{e+1},\ldots$ are all positive, nonzero quantities converging
to $1$, and $e$ denotes a finite integer. The solutions of the homogeneous
equations \eqref{eq:94} are obtained from
\[
  x'_1=u_1,\ldots,x'_{e-1}=u_{e-1},\qquad
  x'_e=0,\quad x'_{e+1}=0,\ldots,
\]
where $u_1,\ldots,u_{e-1}$ are arbitrary constants; hence we see that there
are only finitely many, namely exactly $e-1$, linearly independent solutions
of \eqref{eq:94}.

Thus, under the assumption that $A(x,y)$ is a completely continuous bilinear
form, the system of equations \eqref{eq:93} has all the essential properties
of a system of finitely many equations in equally many unknowns.

In conclusion, it remains to show with what surprising elegance and
simplicity \hyperref[thm:X]{Theorem X} can be proved, without any new
convergence consideration, by making use of \hyperref[thm:V]{Theorems V}
and \hyperref[thm:IX]{IX}. Indeed, from \hyperref[thm:IX]{Theorem IX} we
immediately derive the following fact:

\begin{lemma}[6]\label{lem:6}
If $\varkappa_1,\varkappa_2,\ldots$ is an infinite sequence of positive
quantities converging to $1$, and
\[
  S(x,y)=\sum_{p,q}s_{pq}x_py_q
\]
denotes a completely continuous skew-symmetric form in the infinitely many
variables $x_1,x_2,\ldots,y_1,y_2,\ldots$, then there always exists a
completely continuous bilinear form $T(x,y)$ in the same variables such that
\begin{equation}
  \bigl\{\varkappa(x,\mathord{\cdot})+S(x,\mathord{\cdot})\bigr\}
  \bigl\{(\mathord{\cdot},y)+T(\mathord{\cdot},y)\bigr\}
  =(x,y),
  \tag{106}\label{eq:106}
\end{equation}
where $\varkappa(x)$ denotes the quadratic form
\[
  \varkappa(x)=\varkappa_1x_1^2+\varkappa_2x_2^2+\cdots.
\]
Relation \eqref{eq:106} is equivalent to saying that the system of equations
\begin{equation}
  \begin{aligned}
    \varkappa_1x_1+s_{12}x_2+s_{13}x_3+\cdots&=y_1,\\
    s_{21}x_1+\varkappa_2x_2+s_{23}x_3+\cdots&=y_2,\\
    s_{31}x_1+s_{32}x_2+\varkappa_3x_3+\cdots&=y_3,\\
    &\ \vdots
  \end{aligned}
  \tag{107}\label{eq:107}
\end{equation}
has the solutions
\originalpage{224}
\[
  x_1=y_1+\frac{\partial T(x,y)}{\partial x_1},
  \qquad
  x_2=y_2+\frac{\partial T(x,y)}{\partial x_2},\ldots.
\]
\end{lemma}

\begin{proof}
In $S(x,y)$ we put
\[
  x_1=\frac{1}{\sqrt{\varkappa_1}}x'_1,
  \qquad
  x_2=\frac{1}{\sqrt{\varkappa_2}}x'_2,\ldots,
\]
and
\[
  y_1=\frac{1}{\sqrt{\varkappa_1}}y'_1,
  \qquad
  y_2=\frac{1}{\sqrt{\varkappa_2}}y'_2,\ldots.
\]
We then obtain a completely continuous skew-symmetric form $S'(x',y')$,
while $\varkappa(x)$ passes into $(x',x')$. System \eqref{eq:107} becomes
a system of equations of the following shape:
\begin{equation}
  \begin{aligned}
    x'_1+s'_{12}x'_2+s'_{13}x'_3+\cdots&=y_1^*,\\
    s'_{21}x'_1+x'_2+s'_{23}x'_3+\cdots&=y_2^*,\\
    s'_{31}x'_1+s'_{32}x'_2+x'_3+\cdots&=y_3^*,\\
    &\ \vdots
  \end{aligned}
  \tag{108}\label{eq:108}
\end{equation}
If we now apply to $S'$ the orthogonal transformation of
\hyperref[thm:IX]{Theorem IX}, the system \eqref{eq:108} belonging to $S'$
passes into a system of equations of the following shape:
\begin{align*}
  \xi_1+k_1\xi'_1&=\eta_1,\\
  -k_1\xi_1+\xi'_1&=\eta'_1,\\
  \xi_2+k_2\xi'_2&=\eta_2,\\
  -k_2\xi_2+\xi'_2&=\eta'_2,\\
  &\ \vdots
\end{align*}
As is readily seen, this system has the solutions
\begin{align*}
  \xi_1&=\eta_1+
  \frac{\partial \mathsf T(\xi\xi',\eta\eta')}{\partial\xi_1},\\
  \xi'_1&=\eta'_1+
  \frac{\partial \mathsf T(\xi\xi',\eta\eta')}{\partial\xi'_1},\\
  \xi_2&=\eta_2+
  \frac{\partial \mathsf T(\xi\xi',\eta\eta')}{\partial\xi_2},\\
  \xi'_2&=\eta'_2+
  \frac{\partial \mathsf T(\xi\xi',\eta\eta')}{\partial\xi'_2},\\
  &\ \vdots
\end{align*}
\originalpage{225}
provided that
\begin{align*}
  \mathsf T(\xi\xi',\eta\eta')
  ={}&-\frac{k_1^2}{1+k_1^2}(\xi_1\eta_1+\xi'_1\eta'_1)
      -\frac{k_2^2}{1+k_2^2}(\xi_2\eta_2+\xi'_2\eta'_2)-\cdots\\
    &-\frac{k_1}{1+k_1^2}(\xi_1\eta'_1-\xi'_1\eta_1)
      -\frac{k_2}{1+k_2^2}(\xi_2\eta'_2-\xi'_2\eta_2)-\cdots
\end{align*}
is put. Since the quantities $k_1,k_2,\ldots$ converge to zero, $\mathsf T$ is a
completely continuous form. Returning to the variables
$x'_1,x'_2,\ldots,y'_1,y'_2,\ldots$, and from these to the original variables
$x_1,x_2,\ldots,y_1,y_2,\ldots$, whereby the form $T$ arises from
$\mathsf T$, proves the lemma.
\end{proof}

\begin{proof}[Second proof of Theorem X]
We note that the system of equations \eqref{eq:93} in
\hyperref[thm:X]{Theorem X} retains its form if we apply any orthogonal
transformation to the variables $x_1,x_2,\ldots$ and, at the same time,
orthogonally combine the left-hand sides of those equations in the
corresponding way, since this amounts to a simultaneous orthogonal
transformation of both sequences of variables in $A(x,y)$. For simplicity,
we suppose that such an orthogonal transformation of the bilinear form
$A(x,y)$ has already been carried out so that the completely continuous
quadratic form $A(x,x)$ arising from $A(x,y)$ contains only the squares of
the variables, and thus appears in the shape
\[
  A(x,x)=\alpha_1x_1^2+\alpha_2x_2^2+\cdots,
\]
or
\begin{equation}
  A(x,y)+A(y,x)=2(\alpha_1x_1y_1+\alpha_2x_2y_2+\cdots).
  \tag{109}\label{eq:109}
\end{equation}
Since $\alpha_1,\alpha_2,\ldots$ are quantities converging to zero, there
are certainly only finitely many among them which are $\leq-1$. Let $e$ be
an integer such that
\begin{equation}
  \alpha_{e+p}>-1\qquad(p=1,2,\ldots).
  \tag{110}\label{eq:110}
\end{equation}
We now separate off the first $e$ equations of \eqref{eq:93} in
\hyperref[thm:X]{Theorem X}, and write the remaining equations in the form
\begin{equation}
  \begin{aligned}
    (a_{e+1,e+1}+1)x_{e+1}+a_{e+1,e+2}x_{e+2}+\cdots&=y_{e+1},\\
    a_{e+2,e+1}x_{e+1}+(a_{e+2,e+2}+1)x_{e+2}+\cdots&=y_{e+2},\\
    &\ \vdots
  \end{aligned}
  \tag{111}\label{eq:111}
\end{equation}
where, for brevity, we have put
\begin{equation}
  \begin{aligned}
    y_{e+1}&=a_{e+1}-a_{e+1,1}x_1-\cdots-a_{e+1,e}x_e,\\
    y_{e+2}&=a_{e+2}-a_{e+2,1}x_1-\cdots-a_{e+2,e}x_e,\\
    &\ \vdots
  \end{aligned}
  \tag{112}\label{eq:112}
\end{equation}
\originalpage{226}
These equations \eqref{eq:111}, since by \eqref{eq:109}
\[
  a_{pq}+a_{qp}=0\quad(p\ne q),
  \qquad
  a_{pp}=\alpha_p,
\]
are, in view of \eqref{eq:110}, of the form \eqref{eq:107}, and therefore
permit the application of the lemma just proved. According to that lemma,
there is a completely continuous bilinear form $T(x,y)$ in the variables
$x_{e+1},x_{e+2},\ldots,y_{e+1},y_{e+2},\ldots$ such that equations
\eqref{eq:111} have the solutions
\begin{equation}
  \begin{aligned}
    x_{e+1}&=y_{e+1}+\frac{\partial T}{\partial x_{e+1}},\\
    x_{e+2}&=y_{e+2}+\frac{\partial T}{\partial x_{e+2}},\\
    &\ \vdots
  \end{aligned}
  \tag{113}\label{eq:113}
\end{equation}
If, with due regard to the values \eqref{eq:112} of
$y_{e+1},y_{e+2},\ldots$, we insert these solutions into the first $e$
equations, separated off above, of the given system \eqref{eq:93}, there
arises a system of $e$ equations in the $e$ unknowns $x_1,\ldots,x_e$, as
follows:
\begin{equation}
  \begin{aligned}
    E_{11}x_1+\cdots+E_{1e}x_e&=E_1,\\
    &\ \vdots\\
    E_{e1}x_1+\cdots+E_{ee}x_e&=E_e,
  \end{aligned}
  \tag{114}\label{eq:114}
\end{equation}
where $E_1,\ldots,E_e$ are homogeneous linear forms in
$a_1,a_2,\ldots$, while $E_{11},\ldots,E_{ee}$ are expressed in a known way
in terms of the coefficients of $A(x,y)$.

If these equations have solutions $x_1,\ldots,x_e$, then the values
$x_{e+1},x_{e+2},\ldots$ are calculated from them by means of
\eqref{eq:112} and \eqref{eq:113}, and we thus arrive at the solutions of the
originally given system of equations \eqref{eq:93}. In the other case, the
homogeneous equations
\begin{align*}
  E_{11}x_1+\cdots+E_{1e}x_e&=0,\\
  &\ \vdots\\
  E_{e1}x_1+\cdots+E_{ee}x_e&=0
\end{align*}
can certainly be satisfied by values $x_1,\ldots,x_e$ not all zero. We then
put zero everywhere in place of $a_1,a_2,\ldots$, whereby in fact
\[
  E_1=0,\ldots,E_e=0,
\]
and by means of \eqref{eq:112} and \eqref{eq:113} we arrive at values
$x_{e+1},x_{e+2},\ldots$ which, together with the values
$x_1,\ldots,x_e$ found above, form a
\originalpage{227}
solution system of the homogeneous equations \eqref{eq:94} in
\hyperref[thm:X]{Theorem X}. Thus \hyperref[thm:X]{Theorem X} is completely
proved.
\end{proof}

Since \hyperref[thm:IX]{Theorem IX} on skew-symmetric forms was also proved
above solely with the aid of \hyperref[thm:V]{Theorem V} on the orthogonal
transformation of completely continuous quadratic forms, without any new
convergence consideration, it follows that the theory of equations of the
form \eqref{eq:93}, and hence the theory of the completely continuous
bilinear form in general, can likewise be founded solely upon the theory of
the orthogonal transformation of completely continuous quadratic forms,
without new convergence considerations--a remarkable fact which lends to
the theory of completely continuous forms in infinitely many variables a
wonderful transparency and unity.


\phantomsection
\addcontentsline{toc}{section}{References}
\begin{thebibliography}{9}

\bibitem{Stieltjes1894}\label{bib:Stieltjes1894}
T.-J. Stieltjes,
``Recherches sur les fractions continues,''
\emph{Annales de la Facult\'e des Sciences de Toulouse}, vol.~8 (1894),
J1--J122.
\href{https://www.numdam.org/item/AFST_1894_1_8_4_J1_0/}{Digital edition}.

\bibitem{Hilbert1904Dirichlet}\label{bib:Hilbert1904Dirichlet}
David Hilbert,
``\"Uber das Dirichletsche Prinzip,''
\emph{Mathematische Annalen}, vol.~59 (1904), pp.~161--186.
\href{https://doi.org/10.1007/BF01444753}{doi:10.1007/BF01444753}.

\end{thebibliography}

\end{document}

